\documentclass[11pt]{amsart}
\usepackage{amssymb}
\usepackage{booktabs}
\usepackage[dvipsnames]{xcolor}
\usepackage{microtype}
\usepackage[colorlinks]{hyperref}
\hypersetup{linkcolor=MidnightBlue, citecolor=MidnightBlue, urlcolor=MidnightBlue}

\linespread{1.06}

\usepackage{listings}
\lstdefinelanguage{lean}{
  keywords={theorem,def,structure,abbrev,lemma,noncomputable,where,fun,let,by,Prop,Type,extends,instance,class},
  sensitive=true,
  comment=[l]{--},
  morecomment=[s]{/-}{-/},
  string=[b]",
}
\lstdefinestyle{lean}{
  language=lean,
  basicstyle=\ttfamily\footnotesize,
  keywordstyle=\bfseries\color{MidnightBlue},
  commentstyle=\itshape\color{gray},
  columns=fullflexible,
  keepspaces=true,
  breaklines=true,
  breakatwhitespace=true,
  frame=single,
  rulecolor=\color{gray!45},
  backgroundcolor=\color{gray!6},
  xleftmargin=2mm,
  aboveskip=3.5mm,
  belowskip=3.5mm,
  extendedchars=true,
  literate={¬}{{$\neg$}}1 {¹}{{${}^{1}$}}1 {ö}{{\"o}}1 {Λ}{{$\Lambda$}}1
    {Ω}{{$\Omega$}}1 {α}{{$\alpha$}}1 {β}{{$\beta$}}1 {λ}{{$\lambda$}}1
    {μ}{{$\mu$}}1 {σ}{{$\sigma$}}1 {φ}{{$\varphi$}}1 {ε}{{$\varepsilon$}}1
    {ᵐ}{{${}^{m}$}}1 {‖}{{$\|$}}1 {•}{{$\bullet$}}1
    {⁻¹}{{${}^{-1}$}}2 {⁻}{{${}^{-}$}}1 {₀}{{${}_{0}$}}1 {₁}{{${}_{1}$}}1
    {₂}{{${}_{2}$}}1 {ℂ}{{$\mathbb{C}$}}1 {ℕ}{{$\mathbb{N}$}}1
    {ℝ}{{$\mathbb{R}$}}1 {→}{{$\to$}}1 {↔}{{$\leftrightarrow$}}1
    {∀}{{$\forall$}}1 {∂}{{$\partial$}}1 {∃}{{$\exists$}}1 {∈}{{$\in$}}1
    {∉}{{$\notin$}}1 {∑}{{$\sum$}}1 {∞}{{$\infty$}}1 {∧}{{$\wedge$}}1
    {∫}{{$\int$}}1 {≠}{{$\neq$}}1 {≤}{{$\leq$}}1 {≥}{{$\geq$}}1
    {⊆}{{$\subseteq$}}1 {⟨}{{$\langle$}}1 {⟩}{{$\rangle$}}1
    {𝓝}{{$\mathcal{N}$}}1,
}


\subjclass[2020]{43A99, 43A85, 30H25, 42B15, 42B35, 42C15, 42C40, 28C99}
\keywords{atomic decomposition, Besov space, harmonic analysis, wavelets, multipliers, Haar wavelets, atoms, formal verification, Lean}

\title[Besov-ish spaces: the formalized results]{Besov-ish spaces through atomic decomposition:\\ the results formalized in Lean}

\author[D. Smania]{Daniel Smania}
\address{Departamento de Matem\'atica \\
  Instituto de Ci\^encias Matem\'aticas e de Computa\c{c}\~ao-Universidade de S\~ao Paulo (ICMC/USP) - S\~ao Carlos \\
  Caixa Postal 668 \\ S\~ao Carlos-SP \\ CEP 13560-970 \\ Brazil.}
\email{smania@icmc.usp.br}
\urladdr{\url{https://sites.icmc.usp.br/smania/}}

\newtheorem{theorem}{Theorem}[section]
\newtheorem{corollary}[theorem]{Corollary}
\newtheorem{lemma}[theorem]{Lemma}
\newtheorem{proposition}[theorem]{Proposition}
\theoremstyle{definition}
\newtheorem{definition}[theorem]{Definition}
\newtheorem{remark}[theorem]{Remark}

\newcommand{\SA}{{\mathcal A}}
\newcommand{\SB}{{\mathcal B}}
\newcommand{\SF}{{\mathcal F}}
\newcommand{\SG}{{\mathcal G}}
\newcommand{\SH}{{\mathcal H}}
\newcommand{\SP}{{\mathcal P}}
\newcommand{\SW}{{\mathcal W}}
\newcommand{\lean}[1]{\texttt{#1}}

\hypersetup{colorlinks=true, pdftitle={Besov-ish spaces through atomic decomposition: the results formalized in Lean}, pdfauthor={Daniel Smania}}

\begin{document}

\begin{abstract}
This document describes, without proofs, exactly those results of the paper
\emph{Besov-ish spaces through atomic decomposition} (Analysis \& PDE 15
(2022), no.\ 1; LaTeX source \texttt{besovish-2021-07-19.tex}) that have been
formally verified in the Lean~4 repository \texttt{BesovSpacesGoodGrid},
together with a few results obtained in the formalization that go beyond the
paper. Statements follow the notation of the paper; footnotes record every
deviation from the published statements, and each result is annotated with the
names of the corresponding Lean declarations. The repository compiles with no
\texttt{sorry}, and the main theorems depend only on the standard axioms
\texttt{propext}, \texttt{Classical.choice}, \texttt{Quot.sound}.
\end{abstract}

\maketitle

\setcounter{tocdepth}{2}
\tableofcontents

\section{Introduction}

This is a mathematical description of the current state (June 9, 2026) of the
Lean~4 formalization of \cite{paper}. It was produced from the LaTeX source of
the published paper by the following rules.

\begin{itemize}
\item Results that are not present in the repository were deleted. These are:
  the elementary estimates of paper Section~4 (H\"older-like and convolution
  tricks, inlined in the Lean proofs without standalone statements); the
  H\"older and bounded variation atoms (paper Sections~11B--11C) and the
  corresponding characterizations (paper Propositions~16.2 and~16.3); Dirac's
  approximations (paper Section~17); Corollary~18.6, strongly regular domains
  (paper 18.7--18.8, except for the grid-cell case, which is formalized),
  Pointwise Multipliers~I and~II (paper 18.9--18.10); and all of paper
  Sections~19--21 (quasialgebra, description of $\SB^s_{1,1}$ by indicators,
  left compositions).
\item Statements were updated to match the precision of the formalization;
  every change with mathematical content is recorded in a footnote comparing
  with \texttt{besovish-2021-07-19.tex}.
\item No proofs are included. For each statement the corresponding Lean
  declarations and files are given in a footnote, and the statement is
  followed by a snapshot of the Lean code that formalizes it (the formal
  statement as it appears in the repository, with the proof omitted; an
  ellipsis \texttt{...} marks the omitted proof or, in a few long
  definitions, omitted structural fields).
\end{itemize}

A global change must be singled out: \emph{throughout this document
$p \in [1,\infty)$, $q \in [1,\infty]$ and $u \in [1,\infty]$}. The paper
develops Part~I for the quasi-Banach range $p,q \in (0,\infty]$; the
formalization restricts to the Banach range, so $\rho=\min\{1,p,q\}=1$
everywhere and all the $\rho$-norms of the paper are genuine
norms.\footnote{In Lean the standing hypotheses are \lean{1 <= p},
\lean{p != \(\infty\)}, \lean{Fact (1 <= q)}, \lean{1 <= u}, \lean{0 < s}.}

\section{Notation}

\begin{table}[h]
  \centering
  \caption{Most important notation/symbols/constants}
  \label{tab:table1}
  \begin{tabular}{cc}
    \toprule
    Symbol & Description\\
    \midrule
    $I$ & phase space\\
    $m$ & finite measure in the phase space $I$\\
    $a_P, b_P$ & an atom supported on $P$\\
    $\mathcal{A}$ & a class of atoms \\
    $\mathcal{A}(Q)$ & set of atoms of class $\mathcal{A}$ supported on $Q$ \\
    $\mathcal{A}^{sz}_{s,p}$& class of $(s,p)$-Souza's atoms \\
    $\mathcal{A}^{bs}_{s,\beta,p,q}$ & class of $(s,\beta,p,q)$-Besov's atoms \\
    $\mathcal{P}$ & grid of $I$\\
    $\mathcal{P}^k$ & family of subsets of $I$ at the $k$-th level of $\mathcal{P}$\\
    $\lambda_1 \leq \lambda_2$ & describes geometry of the good grid $\mathcal{P}$\\
    $\mathcal{B}^s_{p,q}(\mathcal{A})$ & $(s,p,q)$-Banach space defined by the class of atoms $\mathcal{A}$\\
    $\mathcal{B}^s_{p,q}$ or $\mathcal{B}^s_{p,q}(\mathcal{A}^{sz}_{s,p})$ & $(s,p,q)$-Banach space defined by Souza's atoms\\
    $P, Q, W$ & elements of the grid $\mathcal{P}$\\
    $L^p$ or $L^p(m)$ & Lebesgue spaces of $(I,\mathbb{A}, m)$ \\
    $|\cdot|_p$ & norm in $L^p$, $p \in (0,\infty]$ \\
    $p'$ & $1/p+1/p'=1$, with $p\in [1,\infty]$ \\
    \bottomrule
  \end{tabular}
\end{table}

We removed from the paper's table the rows for H\"older atoms, bounded
variation atoms, the constants $\rho$ and $\hat t$ (no longer needed in the
Banach range) since the corresponding material is not formalized. Constants
denoted $C$, $C(\cdot)$ below depend only on the displayed arguments; the Lean
statements carry explicit values for most of them.

\vspace{8mm}
\centerline{\Large\bf I. DIVIDE AND RULE}
\vspace{6mm}

\begin{center}
\fbox{\begin{minipage}{0.88\textwidth}
\centering
In Part I we assume $s>0$, \ $p \in [1,\infty)$, \ $u\in[1,\infty]$ \ and \
$q \in [1,\infty]$.\footnote{Paper: $s>0$, $p\in (0,\infty)$,
$q\in(0,\infty]$. See the Introduction for the restriction to the Banach
range.}
\end{minipage}}
\end{center}
\vspace{4mm}

\section{Measure spaces and grids}\label{partition}

Let $I$ be a measure space with a $\sigma$-algebra $\mathbb{A}$ and $m$ be a
measure on $(I,\mathbb{A})$, $m(I)<\infty$. Given a measurable set $J\subset I$
denote $|J|=m(J)$. We denote the Lebesgue spaces of $(I,\mathbb{A},m)$ by
$L^p$. A {\bf grid} is a sequence of finite families of measurable sets with
positive measure $\mathcal{P}=(\mathcal{P}^k)_{k\in \mathbb{N}}$, so that at
least one of these families is non empty and
\begin{itemize}
\item[$G_1$.] Given $Q\in \mathcal{P}^k$, let
$$\Omega_Q^k=\{ P \in \mathcal{P}^k\colon \ P\cap Q\neq \emptyset \}.$$
Then
$$C_{1}=\sup_k \sup_{Q \in \mathcal{P}^k} \# \Omega_Q^k< \infty.$$
\end{itemize}
Define $|\mathcal{P}^k|=\sup \{|Q|\colon Q \in \mathcal{P}^k\}$. We often abuse
notation using $\mathcal{P}$ for both $(\mathcal{P}^k)_{k\in \mathbb{N}}$ and
$\cup_k \mathcal{P}^k$.\footnote{Lean: structure \lean{WeakGrid} (fields
\lean{partitions}, \lean{measurable}, \lean{positive\_measure},
\lean{exists\_nonempty}, \lean{Cmult1}, \lean{overlap\_card\_le}) and
\lean{WeakGridSpace}, in \lean{WeakGrid/Definition.lean}. The paper calls
this simply a ``grid''; in the repository it is called a \emph{weak grid}, to
distinguish it from the good grids of Section~\ref{goodgrids}. The paper's
convention that $P \neq Q$ for cells at different levels is replaced by
indexing cells by (level, cell) pairs (\lean{WeakGridCell}).}
\begin{lstlisting}[style=lean]
structure WeakGrid where
        /-- The measure on `α`. -/
        μ : MeasureTheory.Measure α
        /-- The measure is finite on the whole space. -/
        isFinite : MeasureTheory.IsFiniteMeasure μ
        /-- The finite family `P^k` at each level `k`. -/
        partitions : ℕ → Finset (Set α)
        /-- Every cell is measurable. -/
        measurable : ∀ k Q, Q ∈ partitions k → MeasurableSet Q
        /-- Every cell has positive measure. -/
        positive_measure : ∀ k Q, Q ∈ partitions k → 0 < μ Q
        /-- At least one level is nonempty. -/
        exists_nonempty : ∃ k, (partitions k).Nonempty
        /-- Uniform same-level overlap multiplicity bound. -/
        Cmult1 : ℕ
        /-- `# Ω_Q^k ≤ Cmult1` for every `Q ∈ P^k`. -/
        overlap_card_le :
          ∀ k Q, Q ∈ partitions k → (overlapFinset (partitions k) Q).card ≤ Cmult1

structure WeakGridSpace where
        grid : WeakGrid (α := α)
\end{lstlisting}

\begin{remark}
There are plenty of examples of spaces with grids; the simplest one is
$[0,1)$ with the Lebesgue measure and the dyadic grid
$\mathcal{D}^k=\{[i/2^k,(i+1)/2^k),\ 0\leq i< 2^k\}$, and similarly the dyadic
and $d$-adic grids of $[0,1)^n$. Any measure space with a finite non-atomic
measure can be endowed with a grid made of a nested sequence of partitions
whose elements at the same level have precisely the same measure.
\end{remark}

\section{Atoms}\label{secatoms}

Let $\mathcal{P}$ be a grid. Let $p \in [1,\infty)$, $u \in [1,\infty]$, and
$s>0$. A family of {\bf atoms} associated with $\mathcal{P}$ of type $(s,p,u)$
is an indexed family $\mathcal{A}$ of pairs
$(\mathcal{B}(Q),\mathcal{A}(Q))_{Q\in \cup_k \mathcal{P}^k}$ where
\begin{itemize}
\item[$A_1$.] $\mathcal{B}(Q)$ is a complex linear space of functions
contained in $L^{pu}$.\footnote{Paper $A_1$ asks $\mathcal{B}(Q)$ to be a
complex \emph{Banach} space contained in $L^{pu}$. Lean
(\lean{LocalVectorSpace}, in \lean{WeakGrid/Atoms.lean}) only requires a
complex vector space together with an injective linear map into the functions
$\alpha \to \mathbb{C}$, all whose members belong to $L^{pu}$
(\lean{local\_memLp}); no completeness of the local space is assumed, and it
is never used.}
\item[$A_2$.] If $\phi \in \mathcal{B}(Q)$ then $\phi(x)=0$ for every
$x \not\in Q$.
\item[$A_3$.] $\mathcal{A}(Q)$ is a nonempty convex subset of $\mathcal{B}(Q)$
such that $\phi \in \mathcal{A}(Q)$ implies
$\sigma\phi\in\mathcal{A}(Q)$ for every $\sigma \in \mathbb{C}$ satisfying
$|\sigma|=1$.\footnote{Lean: \lean{atoms\_nonempty}, \lean{atoms\_convex},
\lean{atoms\_phase\_invariant}. The paper states the phase invariance as an
``if and only if''; the two formulations are equivalent.}
\item[$A_4$.] We have $$|\phi|_{pu} \leq |Q|^{s-\frac{1}{u'p}}$$ for every
$\phi \in \mathcal{A}(Q)$, where $1/u + 1/u' = 1$.
\end{itemize}
We will say that $\phi \in \mathcal{A}(Q)$ is an $\mathcal{A}$-atom of type
$(s,p,u)$ supported on $Q$. Sometimes we also assume
\begin{itemize}
\item[$A_5$.] For every $Q\in \mathcal{P}$ we have that $\mathcal{A}(Q)$ is a
compact subset in the strong topology of $L^p$,
\end{itemize}
or
\begin{itemize}
\item[$A_6$.] every $\mathcal{A}(Q)$ is a sequentially compact subset in the
weak topology of $L^p$.\footnote{Lean: structure \lean{AtomFamily} in \lean{WeakGrid/Atoms.lean};
$A_5$ appears as \lean{AssumptionA5}. The paper's alternative assumption
$A_7$ ($\mathcal{B}(Q)$ finite dimensional and $\mathcal{A}(Q)$ closed) is
used in the paper only to derive $A_5$; in the formalization the
finite-dimensional case is instantiated directly for Souza's atoms.}
\end{itemize}
\begin{lstlisting}[style=lean,basicstyle=\ttfamily\scriptsize]
structure AtomFamily
  (G : WeakGridSpace (α := α)) (s : ℝ) (p u : ℝ≥0∞) where
  /-- The Hölder conjugate exponent `u'`. -/
  uConj : ℝ≥0∞
  /-- The smoothness parameter is positive. -/
  s_pos : 0 < s
  /-- `p ∈ [1,∞)`. -/
  one_le_p : 1 ≤ p
  /-- `p < ∞`, expressed in `ℝ≥0∞` as `p ≠ ∞`. -/
  p_ne_top : p ≠ ∞
  /-- `u ∈ [1,∞]`. -/
  one_le_u : 1 ≤ u
  /-- `uConj` is the Hölder conjugate of `u`. -/
  holder_conjugate : ENNReal.HolderConjugate u uConj
  /-- The local vector space `B(Q)`. -/
  localSpace : WeakGridCell G → LocalVectorSpace.{u, v} α
  /-- The chosen atoms `A(Q)`, as elements of `B(Q)`. -/
  atoms : ∀ Q, Set ((localSpace Q).carrier)
  /-- Every cell has at least one atom. -/
  atoms_nonempty : ∀ Q, (atoms Q).Nonempty
  /-- `B(Q)` is contained in `L^{pu}`. -/
  local_memLp : ∀ Q φ, MeasureTheory.MemLp ((localSpace Q).toFun φ) (p * u) G.measure
  /-- Local functions are supported on `Q`. -/
  local_support : ∀ Q φ, ∀ x, x ∉ Q.cell → (localSpace Q).toFun φ x = 0
  /-- `A(Q)` is convex in the local vector space. -/
  atoms_convex : ∀ Q, Convex ℝ (atoms Q)
  /-- `A(Q)` is invariant under multiplication by complex scalars of modulus one. -/
  atoms_phase_invariant :
    ∀ Q φ (σ : ℂ), φ ∈ atoms Q → ‖σ‖ = (1 : ℝ) → σ • φ ∈ atoms Q
  /-- The atom size estimate in `L^{pu}`. -/
  atom_bound :
    ∀ Q φ, φ ∈ atoms Q →
      MeasureTheory.eLpNorm ((localSpace Q).toFun φ) (p * u) G.measure ≤
        atomMeasureScale G s p uConj Q
\end{lstlisting}

\section{Besov-ish spaces}\label{besovish}

Let $p \in [1,\infty)$, $u \in [1,\infty]$, $q \in [1,\infty]$, $s>0$,
$\mathcal{P}=(\mathcal{P}^k)_{k\geq 0}$ be a grid and let $\mathcal{A}$ be a
family of atoms of type $(s,p,u)$. We will also assume
\begin{itemize}
\item[$G_2$.] We have
$$\Big( \sum_k |\mathcal{P}^k|^{s q^\ast}\Big)^{1/q^\ast} < \infty
\quad\text{(usual modification if } q^\ast=\infty\text{)},
\qquad \lim_k |\mathcal{P}^k|=0,$$
where $q^\ast$ is the exponent produced by the H\"older-like estimate of the
paper (for $q>1$, $q^\ast = q'$).\footnote{Lean: \lean{AssumptionG2} in
\lean{WeakGrid/BesovishSpaces.lean}. The paper writes this condition as
$C(p,q,(|\mathcal{P}^k|^s)_k)<\infty$ using the constant of its
Proposition~4.1; since Section~4 has no standalone Lean counterpart, the
formalization states the required summability of the weighted level measures
directly, together with $\lim_k|\mathcal{P}^k| = 0$.}
\end{itemize}

The {\bf Besov-ish space} $\mathcal{B}^s_{p,q}(I,\mathcal{P},\mathcal{A})$ is
the set of all complex valued functions $g \in L^p$ that can be represented by
an absolutely convergent series in $L^p$
\begin{equation} \label{rep} g = \sum_{k=0}^{\infty} \sum_{Q \in \mathcal{P}^k} s_Q a_Q\end{equation}
where $a_Q$ is in $\mathcal{A}(Q)$, $s_Q \in \mathbb{C}$ and with finite
{\bf cost}
\begin{equation} \label{rep2} \big( \sum_{k=0}^{\infty} (\sum_{Q \in \mathcal{P}^k} |s_Q|^p)^{q/p} \big)^{1/q} < \infty\end{equation}
(usual modification when $q = \infty$). The series in (\ref{rep}) is called a
{\bf $\mathcal{B}^s_{p,q}(I,\mathcal{P},\mathcal{A})$-representation} of the
function $g$. Define
$$|g|_{\mathcal{B}^s_{p,q}(I,\mathcal{P},\mathcal{A})} = \inf \big( \sum_{k=0}^{\infty} (\sum_{Q \in \mathcal{P}^k} |s_Q|^p )^{q/p} \big)^{1/q},$$
where the infimum runs over all possible representations of $g$ as in
(\ref{rep}).\footnote{Lean: \lean{LpGridRepresentation} (a representation is a
sequence of level blocks together with a \lean{HasSum} witness in $L^p$),
\lean{pqCost}, \lean{MemBesovishCoeffCost}, \lean{BesovishSpace} (a
\lean{Submodule} of $L^p$) and the cost norm \lean{Norm\_Costpq} (an infimum
over representation costs), all in \lean{WeakGrid/BesovishSpaces.lean}.}
When the choice of $I$ and/or $\mathcal{P}$ is obvious we write
$\mathcal{B}^s_{p,q}(\mathcal{P},\mathcal{A})$ or
$\mathcal{B}^s_{p,q}(\mathcal{A})$; whenever we write just
$\mathcal{B}^s_{p,q}$ we choose the Souza's atoms $\mathcal{A}^{sz}_{s,p}$ of
Section~\ref{secatom}.
\begin{lstlisting}[style=lean]
structure LpGridRepresentation
    (A : AtomFamily G s p u) (g : Lp ℂ p G.measure) where
  block : (k : ℕ) → LevelBlock A k
  hasSum : HasSum (fun k => (block k).toLp A) g

def BesovishSpace (A : AtomFamily G s p u) (q : ℝ≥0∞)
    [Fact (1 ≤ q)]
    : Submodule ℂ (Lp ℂ p G.measure) where
  -- Carrier: all `L^p` elements admitting a Besov-ish atomic representation with finite p q cost.
  carrier := { g | MemBesovishCoeffCost A q g }
  ...

noncomputable def Norm_Costpq
    (A : AtomFamily G s p u) (q : ℝ≥0∞) [Fact (1 ≤ q)]
    (g : BesovishSpace A q) : ℝ :=
  pqPseudoNorm A q g
\end{lstlisting}

\begin{proposition}[Embedding into $L^t$]\label{lp}
Assume $G_1$--$G_2$ and $A_1$--$A_4$. Let $t \in [p,\infty)$ be such that
\begin{equation}\label{c11} s-\frac{1}{p}+\frac{1}{t} \geq 0, \qquad p \leq t \leq pu,\end{equation}
and suppose
\begin{equation}\label{decay} C_{2}(t)= C_1^{1+1/t}\, C\big(t,q,(|\mathcal{P}^k|^{t(s-\frac1p+\frac1t)})_k\big) < \infty.\end{equation}
Then for every $g$ with a $\mathcal{B}^s_{p,q}(\mathcal{A})$-representation
(\ref{rep}) of cost $K$ the series converges in $L^t$ and
\begin{equation}\label{inclulp} |g|_{t} \leq C_2(t)\, K.\end{equation}
In particular
$$|g|_t \leq C_2(t)\, |g|_{\mathcal{B}^s_{p,q}(\mathcal{A})},$$
and $\mathcal{B}^s_{p,q}(\mathcal{A}) \subset L^t$
continuously.\footnote{Lean: \lean{lp\_embedding\_adapted\_statement} in
\lean{WeakGrid/BesovishSpaces.lean}. The Lean statement is the
per-representation inequality (\ref{inclulp}) with the explicit constant
$C_1^{1+1/t}\cdot$\lean{cCoefficient}; the norm inequality follows by taking
the infimum over representations. Compared with paper Proposition~6.1 the
hypotheses are identical except for $q \geq 1$ and $t \geq p$ being imposed
globally, and the levelwise estimate (6.2) of the paper is internal to the
proof rather than part of the statement.}
\end{proposition}
\begin{lstlisting}[style=lean]
theorem lp_embedding_adapted_statement
    {A : AtomFamily G s p u} {t : ℝ≥0∞}
    [Fact (1 ≤ t)]
    (hp_ne_top : p ≠ ∞) (ht_ne_top : t ≠ ∞)
    (hq_one : 1 ≤ q)
    (hp_le_t : p ≤ t) (ht_le_pu : t ≤ p * u)
    (hs_nonneg : 0 ≤ s - 1 / p.toReal + 1 / t.toReal)
    {g : Lp ℂ p G.measure} (R : LpGridRepresentation A g)
    (hRfin : LpGridRepresentation.FinitePQCost (q := q) R)
    (hCco_fin : cCoefficientFinite t q (fun k =>
      (levelMeasureWeight G s p t k) ^ t.toReal)) :
        (MeasureTheory.eLpNorm (g : α → ℂ) t G.measure).toReal ≤
        ((G.grid.Cmult1 : ℝ) ^ (1 + 1 / t.toReal)) *
        cCoefficient t q (fun k => (levelMeasureWeight G s p t k) ^ t.toReal) *
          LpGridRepresentation.pqCost (q := q) R := ...
\end{lstlisting}

\begin{proposition}[Linear structure]\label{linear}
Assume $G_1$--$G_2$ and $A_1$--$A_4$. Then
$\mathcal{B}^s_{p,q}(\mathcal{A})$ is a complex linear subspace of $L^p$ and
$|\cdot|_{\mathcal{B}^s_{p,q}(\mathcal{A})}$ is a norm on it. Moreover the
inclusion
$$\imath\colon (\mathcal{B}^s_{p,q}(\mathcal{A}),|\cdot|_{\mathcal{B}^s_{p,q}(\mathcal{A})})\rightarrow (L^p,|\cdot|_p)$$
is continuous.\footnote{Lean: \lean{BesovishSpace} is a \lean{Submodule} of
$L^p$; the normed structure is \lean{BesovishSpace.costNormedAddCommGroup}
(requiring $p \neq \infty$ and $G_2$), and continuity of $\imath$ is the case
$t=p$ of Proposition~\ref{lp}. Compared with paper Proposition~6.3,
$\rho=\min\{1,p,q\}=1$ since $p,q\geq 1$, so the $\rho$-norm of the paper is a
norm.}
\end{proposition}
\begin{lstlisting}[style=lean]
noncomputable def costNormedAddCommGroup
    (hp_top : p ≠ ∞)
    (hCco_fin : LpGridRepresentation.cCoefficientFinite p q (fun k =>
      (LpGridRepresentation.levelMeasureWeight G s p p k) ^ p.toReal)) :
    NormedAddCommGroup (BesovishSpace A q) where
  norm := Norm_Costpq A q
  dist x y := Norm_Costpq A q (-x + y)
  ...
\end{lstlisting}

\begin{proposition}[Limits of representations]\label{compa2}
Assume $G_1$--$G_2$ and $A_1$--$A_4$. Suppose that $g_n$ are functions in
$\mathcal{B}^s_{p,q}(\mathcal{A})$ with
$\mathcal{B}^s_{p,q}(\mathcal{A})$-representations
$$g_n=\sum_{k=0}^{\infty}\sum_{Q\in \SP^k}s_Q^na_Q^n$$
satisfying
\begin{itemize}
\item[i.] There is $C$ such that for every $n$
$$\big(\sum_{k=0}^{\infty}(\sum_{Q\in \SP^k}|s_Q^n|^p)^{q/p}\big)^{1/q}\leq C.$$
\item[ii.] For every $Q\in \mathcal{P}$ the limit $s_Q=\lim_n s_Q^n$ exists.
\item[iii.] For every $Q\in \mathcal{P}$ there is $a_Q\in \mathcal{A}(Q)$ such
that either $a_Q^n \to a_Q$ in the strong topology of $L^p$, or $a_Q^n \to
a_Q$ weakly in $L^p$.
\end{itemize}
Then $g_n$ converges (strongly, respectively weakly) in $L^p$ to a function
$g \in \mathcal{B}^s_{p,q}(\mathcal{A})$ which has the
$\mathcal{B}^s_{p,q}(\mathcal{A})$-representation
$$g=\sum_{k=0}^{\infty}\sum_{Q\in \SP^k}s_Q\, a_Q$$
of cost at most $C$.\footnote{Lean: \lean{representation\_limit} and
\lean{representation\_limit\_strong} in \lean{WeakGrid/Completeness.lean},
matching paper Proposition~6.4.}
\end{proposition}
\begin{lstlisting}[style=lean]
theorem representation_limit
    (A : AtomFamily G s p u)(hG2 : AssumptionG2 G s p u q)
     {gseq : ℕ → Lp ℂ p G.measure}
    {gLim : Lp ℂ p G.measure} {C : ℝ}
    (H : RepresentationLimitHypotheses A q gseq gLim C)
    (hp_ne_top : p ≠ ∞) (hs_pos : 0 < s) (hu_one : 1 ≤ u)
    [Fact (1 ≤ u)] (hC : 0 ≤ C) :
    MemBesovishCoeffCost A q gLim ∧
      LpGridRepresentation.FinitePQCost (q := q) H.Rlim ∧
      LpGridRepresentation.pqCost (q := q) H.Rlim ≤ C ∧
      Tendsto (fun n => toWeakSpace ℂ (Lp ℂ p G.measure) (gseq n)) atTop
        (𝓝 (toWeakSpace ℂ (Lp ℂ p G.measure) gLim)) := ...
\end{lstlisting}

\begin{corollary}[Compactness and completeness]\label{compa1}
Assume $G_1$--$G_2$, $A_1$--$A_4$, and either $A_5$ or $A_6$. Then
\begin{itemize}
\item[i.] If $g_n \in \mathcal{B}^s_{p,q}(\mathcal{A})$ satisfy
$|g_n|_{\mathcal{B}^s_{p,q}(\mathcal{A})}\leq C$ for every $n$, then there is
a subsequence converging (strongly under $A_5$, weakly under $A_6$) in $L^p$
to some $g \in \mathcal{B}^s_{p,q}(\mathcal{A})$ with
$|g|_{\mathcal{B}^s_{p,q}(\mathcal{A})}\leq C$.\footnote{Lean:
\lean{representation\_limit\_strong\_existence} and
\lean{representation\_limit\_weak\_existence} in
\lean{WeakGrid/Completeness.lean} (paper Corollary~6.5).}
\item[ii.] Under $A_5$, every closed ball of
$(\mathcal{B}^s_{p,q}(\mathcal{A}),|\cdot|_{\mathcal{B}^s_{p,q}(\mathcal{A})})$
is sequentially compact in the strong topology of
$L^p$.\footnote{Lean: \lean{exists\_strongly\_convergent\_subseq\_of\_uniform\_pqCost}
and \lean{closed\_Norm\_Costpq\_ball\_strongly\_seqCompact} in
\lean{WeakGrid/Completeness.lean} (paper Corollary~6.6; the paper states
i.\ and the compactness of the inclusion $\imath$, the formalization states
the sequential compactness of closed cost balls, which is the form used
downstream).}
\item[iii.] Under $A_5$,
$(\mathcal{B}^s_{p,q}(\mathcal{A}),|\cdot|_{\mathcal{B}^s_{p,q}(\mathcal{A})})$
is a complex Banach space.\footnote{Lean:
\lean{besovishSpace\_costNorm\_completeSpace} in
\lean{WeakGrid/Completeness.lean} (paper Corollary~6.7; for $p,q \geq 1$ the
$\rho$-Banach space of the paper is a Banach space). The paper's
Corollary~6.6 (existence of a representation attaining the infimum cost) is
not formalized and is omitted.}
\end{itemize}
\end{corollary}
\begin{lstlisting}[style=lean]
theorem closed_Norm_Costpq_ball_strongly_seqCompact
    (hp_ne_top : p ≠ ∞) (hs_pos : 0 < s) (hu_one : 1 ≤ u)
    (A : AtomFamily G s p u) (hG2 : AssumptionG2 G s p u q)
    (hA5 : AssumptionA5 A)
    {C : ℝ}
    [Fact (1 ≤ u)]
    (hC : 0 ≤ C)
    (gseq : ℕ → BesovishSpace A q)
    (hball : ∀ n, BesovishSpace.Norm_Costpq A q (gseq n) ≤ C) :
    ∃ (φ : ℕ → ℕ) (_hφ : StrictMono φ) (gLim : BesovishSpace A q),
      BesovishSpace.Norm_Costpq A q gLim ≤ C ∧
        Tendsto
          (fun n => ((gseq (φ n) : BesovishSpace A q) : Lp ℂ p G.measure))
          atTop
          (𝓝 ((gLim : BesovishSpace A q) : Lp ℂ p G.measure)) := ...

theorem besovishSpace_costNorm_completeSpace
    (hp_ne_top : p ≠ ∞) (hs_pos : 0 < s) (hu_one : 1 ≤ u)
    (A : AtomFamily G s p u) (hG2 : AssumptionG2 G s p u q)
    (hA5 : AssumptionA5 A)
    [Fact (1 ≤ u)] :
    @CompleteSpace (BesovishSpace A q)
      (BesovishSpace.costNormedAddCommGroup
        (A := A) (q := q) hp_ne_top hG2.1).toMetricSpace.toPseudoMetricSpace.toUniformSpace := ...
\end{lstlisting}

\section{Scales of spaces}\label{escala}

A family of atoms $\mathcal{A}$ of type $(s,p,u)$ induces a one-parameter
scale $\tilde{s}\mapsto \mathcal{A}_{\tilde{s},p}$, where
$\mathcal{A}_{\tilde{s},p}$ is the family of atoms of type $(\tilde{s},p,u)$
defined by
$$\mathcal{A}_{\tilde{s},p}(Q)= \{ |Q|^{\tilde{s}-s}a_Q\colon \ a_Q\in \mathcal{A}(Q)\}.$$

\begin{proposition}[Scale inclusion]\label{compa}
Assume $G_1$--$G_2$ and that $\mathcal{A}$ satisfies $A_1$--$A_4$. Let
$0< s \leq \tilde{s}$ and $q,\tilde{q} \in [1,\infty]$. Suppose
$$\Big( \sum_{k} |\mathcal{P}^k|^{q^\ast(\tilde{s}-s)} \Big)^{1/q^\ast} < \infty.$$
Then $\mathcal{B}^{\tilde{s}}_{p,\tilde{q}}(\mathcal{A}_{\tilde{s},p})\subset
\SB^s_{p,q}(\mathcal{A}_{s,p})$, the inclusion is a continuous linear
map,\footnotemark\ and
$$|g|_{\SB^s_{p,q}(\mathcal{A}_{s,p})} \leq
C\big(q,(|\mathcal{P}^k|^{\tilde s - s})_k\big)\,
|g|_{\mathcal{B}^{\tilde{s}}_{p,\tilde{q}}(\mathcal{A}_{\tilde{s},p})}.$$
\footnotetext{Lean: \lean{smoothnessScaleBesovishSpace\_subset} and
\lean{smoothnessScaleBesovishSpaceInclusion\_Norm\_Costpq\_le} in
\lean{WeakGrid/Scales.lean}, with the explicit constant
\lean{scaleCoefficient}. This is item~A of paper Proposition~7.1; items~B
(compactness of bounded sets of the finer space in the coarser norm) and C
(compactness of the inclusion under $A_7$) are not formalized and are
omitted. The paper's two-parameter scale $(\tilde s,\tilde p)$ is also not
formalized.}
\end{proposition}
\begin{lstlisting}[style=lean]
theorem smoothnessScaleBesovishSpace_subset
  [Fact (1 ≤ p)]
  [Fact (1 ≤ u)] [Fact (1 ≤ q)] [Fact (1 ≤ qTilde)]
    (A : AtomFamily G s p u) (sTilde : ℝ) (hsTilde_pos : 0 < sTilde)
    (hs_le : s ≤ sTilde)
    (hqfin : scaleCoefficientFinite q (smoothnessScaleLevelWeight G s sTilde p))
    {g : Lp ℂ p G.measure}
    (hg : g ∈ SmoothnessScaleBesovishSpace
      (A.smoothnessScaleAtomFamily sTilde hsTilde_pos) qTilde) :
    g ∈ BesovishSpace A q := ...
\end{lstlisting}

\section{Transmutation of atoms}\label{transsec}

The key result in Part~I states that if we replace each atom of an atomic
decomposition by a combination of atoms of a possibly different class, in
such way that we do not change the location of the supports and the mass of
the new representation concentrates pretty much on the original scale, then
we obtain a new atomic decomposition whose cost is controlled by the cost of
the original one.

\begin{proposition}[Transmutation of atoms]\label{trans}
Assume
\begin{itemize}
\item[{\bf I.}] $\mathcal{A}_2$ is a class of $(s,p,u_2)$-atoms for a grid
$\mathcal{W}$, satisfying $A_1$--$A_4$ and $G_1$--$G_2$. Let $\mathcal{G}$ be
also a grid satisfying $G_1$--$G_2$.
\item[{\bf II.}] $k_i\in \mathbb{N}$, $i\geq 0$, is a sequence such that there
are $\alpha>0$ and $A,B\in \mathbb{R}$ with
$$ \alpha i +A \leq k_i\leq \alpha i +B \quad\text{for every } i.$$
\item[{\bf III.}] There is $\lambda \in(0,1)$ such that the following holds.
For every $Q\in \mathcal{G}^i$ there are atoms $b_{P,Q} \in \mathcal{A}_2(P)$
and coefficients $s_{P,Q}\in \mathbb{C}$, indexed by the cells
$P \in \mathcal{W}$ with $P\subset Q$, such that
$$h_Q=\sum_k \sum_{P\in \mathcal{W}^k,\, P\subset Q} s_{P,Q} b_{P,Q}$$
is a $\mathcal{B}^s_{p,q}(\mathcal{A}_2)$-representation of a function $h_Q$,
with $s_{P,Q}=0$ whenever $P \in \mathcal{W}^k$ with $k < k_i$, and moreover
\begin{equation}\label{ad} \sum_{P\in \mathcal{W}^k,\, P \subset Q} |s_{P,Q}|^p \leq C_3\, \lambda^{k-k_i}
\quad\text{for every } k\geq k_i.\end{equation}
\end{itemize}
Then
\begin{itemize}
\item[{\bf A.}] For all coefficients $(c_Q)_{Q\in \mathcal{G}}$ with
$$\big( \sum_i \big( \sum_{Q\in \mathcal{G}^i} |c_{Q}|^p \big)^{q/p}\big)^{1/q}<\infty,$$
the sequence $N\mapsto \sum_{i\leq N} \sum_{Q\in \mathcal{G}^i }c_Q h_Q$
converges in $L^p$ to a function in $\mathcal{B}^s_{p,q}(\mathcal{A}_2)$ that
has a $\mathcal{B}^s_{p,q}(\mathcal{A}_2)$-representation
$\sum_k \sum_{P}m_P d_P$ with $m_P \geq 0$ and
$$\big(\sum_k \big( \sum_{P \in \mathcal{W}^k} |m_{P}|^p \big)^{q/p}\big)^{1/q}
\leq C_4\, \big( \sum_i \big( \sum_{Q\in \mathcal{G}^i} |c_{Q}|^p \big)^{q/p}\big)^{1/q},$$
where $C_4=C_4(C_1,C_3,\lambda,\alpha,A,B,p,q)$ is
explicit.\footnote{Lean: \lean{Transmutation\_of\_Atoms\_Claim\_A} in
\lean{WeakGrid/Transmutation.lean}, with the explicit constant of the paper
(product of $C_1$, $C_3^{1/p}$, $\lambda^{-B/p}$, the convolution-trick
coefficient and the integer-part factor $\max\{\ell\in\mathbb N : \ell <
\alpha\}+1$).}
\item[{\bf B.}] Suppose, in addition to the assumptions of A., that the
$s_{P,Q}$ are nonnegative real numbers and $b_{P,Q}>0$ on $P$ for all $P,Q$.
If $m_P\neq 0$ and $d_P \neq 0$ somewhere on $P$, then $P\subset Q$ for some
cell $Q \in \mathcal{G}$ with $c_Q \neq 0$ and $s_{P,Q}>0$, and the
corresponding $h_Q$ is a positive multiple of a positive function a.e.\ on
$P$. If moreover $c_Q \geq 0$ for every $Q$, then $m_P\neq 0$ implies
$d_P > 0$ a.e.\ on $P$.\footnote{Lean:
\lean{Transmutation\_of\_Atoms\_Claim\_B} and
\lean{Transmutation\_of\_Atoms\_Claim\_B\_sharp} in
\lean{WeakGrid/Transmutation.lean}. Compared with paper Proposition~8.1.B,
the pointwise positivity conclusions are stated almost everywhere on $P$
(with positivity witnessed by positive real multiples), which is the natural
formulation in $L^p$.}
\item[{\bf C.}] Let $\mathcal{A}_1$ be a class of $(s,p,u_1)$-atoms for the
grid $\mathcal{G}$ satisfying $A_1$--$A_4$. Suppose that for every atom
$a_Q\in \mathcal{A}_1(Q)$ we can find $s_{P,Q}$ and $b_{P,Q}$ as in III.\ such
that $h_Q=a_Q$. Then
$$\mathcal{B}^s_{p,q}(\mathcal{A}_1)\subset \mathcal{B}^s_{p,q}(\mathcal{A}_2)$$
and this inclusion is continuous, with
$|\phi|_{\mathcal{B}^s_{p,q}(\mathcal{A}_2)} \leq C_4\,
|\phi|_{\mathcal{B}^s_{p,q}(\mathcal{A}_1)}$ for every $\phi \in
\mathcal{B}^s_{p,q}(\mathcal{A}_1)$.\footnote{Lean:
\lean{Transmutation\_of\_Atoms\_Claim\_C\_explicit} and
\lean{Transmutation\_of\_Atoms\_continuous\_embedding\_explicit} in
\lean{WeakGrid/Transmutation.lean}.}
\end{itemize}
\end{proposition}
\begin{lstlisting}[style=lean,basicstyle=\ttfamily\scriptsize]
theorem Transmutation_of_Atoms_Claim_A
    (G W : WeakGridSpace (α := α))
    (AW : AtomFamily W s p u)
    (k : ℕ → ℕ)
    (A_als B_als r_als : ℝ)
    (hr_als : 0 < r_als)
    (hk_bound : ∀ i : ℕ,
      (k i : NNReal) ≤ r_als * (i : NNReal) + B_als ∧
      r_als * (i : NNReal) + A_als ≤ (k i : NNReal))
    (lam : ℝ) (hlam_pos : 0 < lam) (hlam_lt : lam < 1)
    (C : ℝ) (hC : 0 ≤ C)
    (h : (i : ℕ) → LevelCell G i → Lp ℂ p W.measure)
    (R : (i : ℕ) → (Q : LevelCell G i) → LpGridRepresentation AW (h i Q))
    (hR : RepresentationWsubGandALS (p := p) (q := q) G W AW k
      ⟨A_als, B_als, r_als, hr_als, hk_bound⟩ lam hlam_pos hlam_lt C hC h R)
    (c : (i : ℕ) → LevelCell G i → ℂ)
    (hc : CoeffFinitePQCost (p := p) (q := q) G c)
    (hq_ne_top : q ≠ ∞)
    (hG2_W : AssumptionG2 W s p u q)
    (hp_ne_top : p ≠ ∞)
    (hs_pos : 0 < s) :
    ∃ gLim : Lp ℂ p W.measure,
      ∃ hsum :
        HasSum (fun j => (TransmutationBlockLimit G W AW h R c A_als r_als j).toLp AW) gLim,
      MemBesovishCoeffCost AW q gLim ∧
      LpGridRepresentation.FinitePQCost (q := q)
        ({ block := TransmutationBlockLimit G W AW h R c A_als r_als
           hasSum := hsum } : LpGridRepresentation AW gLim) ∧
      Tendsto (fun N => PartialSumLevels G W h c N) atTop (𝓝 gLim) ∧
      LpGridRepresentation.pqCost (q := q)
        ({ block := TransmutationBlockLimit G W AW h R c A_als r_als
           hasSum := hsum } : LpGridRepresentation AW gLim) ≤
        (G.grid.Cmult1 : ℝ) *
        C ^ (1 / p.toReal) *
        lam ^ (-(B_als : ℝ) / p.toReal) *
        LpGridRepresentation.cCoefficientInt p ∞
          (transmutationKernelZ lam A_als r_als) *
        (Nat.ceil (r_als : ℝ) : ℝ) ^ (1 / q.toReal) *
        CoeffPQCost (p := p) (q := q) G c := ...

theorem Transmutation_of_Atoms_Claim_B (G W : WeakGridSpace (α := α))
    (AW : AtomFamily W s p u)
    (k : ℕ → ℕ)
    (A_als B_als r_als : ℝ)
    (hr_als : 0 < r_als)
    (hk_bound : ∀ i : ℕ,
      (k i : NNReal) ≤ r_als * (i : NNReal) + B_als ∧
      r_als * (i : NNReal) + A_als ≤ (k i : NNReal))
    (lam : ℝ) (hlam_pos : 0 < lam) (hlam_lt : lam < 1)
    (C : ℝ) (hC : 0 ≤ C)
    (h : (i : ℕ) → LevelCell G i → Lp ℂ p W.measure)
    (R : (i : ℕ) → (Q : LevelCell G i) → LpGridRepresentation AW (h i Q))
    (hR : RepresentationWsubGandALS_pos (p := p) (q := q) G W AW k
      ⟨A_als, B_als, r_als, hr_als, hk_bound⟩ lam hlam_pos hlam_lt C hC h R)
    (c : (i : ℕ) → LevelCell G i → ℂ)
    (hc : CoeffFinitePQCost (p := p) (q := q) G c)
    (hG2_W : AssumptionG2 W s p u q)
    (hp_ne_top : p ≠ ∞)
    (hs_pos : 0 < s) :
    (∃ gLim : Lp ℂ p W.measure,
      HasSum (fun j => (TransmutationBlockLimit G W AW h R c A_als r_als j).toLp AW) gLim ∧
      MemBesovishCoeffCost AW q gLim ∧
      Tendsto (fun N => PartialSumLevels G W h c N) atTop (𝓝 gLim)) ∧
    (∀ j : ℕ, ∀ P : LevelCell W j,
      TransmutationCoeffLimit G W AW h R c A_als r_als P ≠ 0 →
      (∃ x, x ∈ P.1 →
        AW.toFunction (levelCellToWeakGridCell W j P)
          (TransmutationAtomLocalLimit G W AW h R c A_als r_als P) x ≠ 0) →
        ∃ i ∈ Finset.range (transmutationStabilizationIndex A_als r_als j),
          ∃ Q ∈ (G.grid.partitions i).attach,
          P.1 ⊆ Q.1 ∧ c i Q ≠ 0 ∧
            ∃ r : NNReal, 0 < r ∧ ((R i Q).block j).coeff P = (r : ℂ) ∧
              ∀ᵐ x ∂ W.measure.restrict P.1,
                ∃ r : NNReal, 0 < r ∧ h i Q x = (r : ℂ)) := ...

theorem Transmutation_of_Atoms_continuous_embedding_explicit
    (G : WeakGridSpace (α := α))
    (u1 u2 : ℝ≥0∞)
    [Fact (1 ≤ u2)]
    (AG1 : AtomFamily G s p u1)
    (AG2 : AtomFamily G s p u2)
    (lam : ℝ) (hlam_pos : 0 < lam) (hlam_lt : lam < 1)
    (C : ℝ) (hC : 0 ≤ C)
    (hAG1_to_AG2 : ∀ i : ℕ, ∀ Q : LevelCell G i,
      ∀ g : Lp ℂ p G.measure,
        (∃ φ : (AG1.localSpace (levelCellToWeakGridCell G i Q)).carrier,
          AG1.IsAtom (levelCellToWeakGridCell G i Q) φ ∧
            g = atomLp AG1 (levelCellToWeakGridCell G i Q) φ) →
          ∃ Rg : LpGridRepresentation AG2 g,
            CoeffFinitePQCost (p := p) (q := q) G
              (fun j S => (Rg.block j).coeff S) ∧
            (∀ j : ℕ, ∀ S : LevelCell G j,
              (¬ S.1 ⊆ Q.1 → (Rg.block j).coeff S = 0) ∧
              (j < i → (Rg.block j).coeff S = 0)) ∧
            ∀ j : ℕ, i ≤ j → Rg.levelCoeffPower j ≤ C * lam ^ (j - i))
    (hG2_G : AssumptionG2 G s p u2 q)
    (hp_ne_top : p ≠ ∞)
    (hs_pos : 0 < s) :
    0 ≤ transmutationClaimCEmbeddingConstant G p q lam C ∧
      ∀ g : BesovishSpace AG1 q,
        ∃ hg2 : MemBesovishCoeffCost AG2 q (g : Lp ℂ p G.measure),
          BesovishSpace.Norm_Costpq AG2 q
              (⟨(g : Lp ℂ p G.measure), hg2⟩ : BesovishSpace AG2 q) ≤
            transmutationClaimCEmbeddingConstant G p q lam C *
              BesovishSpace.Norm_Costpq AG1 q g := ...
\end{lstlisting}

\section{Good grids}\label{goodgrids}

A $(\lambda_1,\lambda_2)$-{\bf good grid}, with $0<\lambda_1\leq
\lambda_2<1$, is a grid $\mathcal{P}=(\mathcal{P}^k)_{k\in \mathbb{N}}$ with
the following properties:
\begin{itemize}
\item[$G_3$.] $\mathcal{P}^0=\{I\}$.
\item[$G_4$.] $I=\cup_{Q \in \mathcal{P}^k} Q$ up to a set of zero
$m$-measure.
\item[$G_5$.] The elements of $\mathcal{P}^k$ are pairwise disjoint.
\item[$G_6$.] For every $Q\in \mathcal{P}^k$, $k>0$, there exists
$P \in \mathcal{P}^{k-1}$ such that $Q\subset P$.
\item[$G_7$.] We have
$$\lambda_1 \leq \frac{|Q|}{|P|} \leq \lambda_2$$
for every $Q\subset P$ with $Q\in \mathcal{P}^{k+1}$ and
$P\in \mathcal{P}^{k}$, $k\geq 0$.
\item[$G_8$.] The family $\cup_k \mathcal{P}^k$ generates the
$\sigma$-algebra $\mathbb{A}$.\footnote{Lean: structure \lean{GoodGrid} in \lean{GoodGrid/Definition.lean},
which extends the structure \lean{UnbalancedHaarWavelet.Grid} of the external
repository (where the nested-partition properties $G_3$--$G_6$, $G_8$ live)
with the ratio bounds $G_7$ (fields \lean{ratio\_lower},
\lean{ratio\_upper}). The paper requires $\lambda_1 < \lambda_2$; the
formalization allows $\lambda_1 \le \lambda_2$.}
\end{itemize}
\begin{lstlisting}[style=lean]
structure GoodGrid extends UnbalancedHaarWavelet.Grid (α := α) where
        lambda1 : ℝ
        /-- Upper ratio bound (strictly less than 1). -/
        lambda2 : ℝ
        /-- `λ₁` is strictly positive. -/
        hlambda1_pos : 0 < lambda1
        /-- `λ₂` is strictly less than `1`. -/
        hlambda2_lt_one : lambda2 < 1
        /-- `λ₁ ≤ λ₂`. -/
        hlambda1_le_lambda2 : lambda1 ≤ lambda2
        /-- Each child cell's measure is at least `λ₁` times its parent's measure. -/
        ratio_lower : ∀ (n : ℕ) (s t : Set α),
                        s ∈ grid.partitions (n + 1) → t ∈ grid.partitions n → s ⊆ t →
                        ENNReal.ofReal lambda1 * μ t ≤ μ s
        /-- Each child cell's measure is at most `λ₂` times its parent's measure. -/
        ratio_upper : ∀ (n : ℕ) (s t : Set α),
                        s ∈ grid.partitions (n + 1) → t ∈ grid.partitions n → s ⊆ t →
                        μ s ≤ ENNReal.ofReal lambda2 * μ t
\end{lstlisting}

\section{Induced spaces}\label{inducedsec}

Consider a Besov-ish space $\mathcal{B}^s_{p,q}(I,\mathcal{P},\mathcal{A})$,
where $\mathcal{P}$ is a good grid. Given $Q\in \mathcal{P}^{k_0}$, consider
the sequence of finite families of subsets
$\mathcal{P}_Q=(\mathcal{P}_Q^i)_{i\geq 0}$ of $Q$ given by
$$\mathcal{P}^i_Q=\{ P\in \mathcal{P}^{k_0+i},\ P\subset Q\},$$
and let $\mathcal{A}_Q$ be the restriction of $\mathcal{A}$ to indices in
$\mathcal{P}_Q$. Then we can consider the {\bf induced} Besov-ish space
$\mathcal{B}^s_{p,q}(Q,\mathcal{P}_Q,\mathcal{A}_Q)$, and the inclusion
$$i\colon \mathcal{B}^s_{p,q}(Q,\mathcal{P}_Q,\mathcal{A}_Q) \rightarrow \mathcal{B}^s_{p,q}(I,\mathcal{P},\mathcal{A})$$
is well defined and a weak contraction:\footnote{Lean: \lean{inducedWeakGrid}, \lean{inducedWeakGridSpace},
\lean{inducedAtomFamily} and \lean{inducedRepresentationToAmbient} in
\lean{WeakGrid/InducedGrid.lean}. The formalization develops the induced
grid, induced atoms, the transfer of representations from the induced grid to
the ambient one with control of cost, and (for Souza's atoms) the transfer of
positivity; this is the form needed in Part~III. The reverse restriction map
is treated, for cells of the grid, in Lemma~\ref{incint3} below.}
$$|f|_{\mathcal{B}^s_{p,q}(I,\mathcal{P},\mathcal{A})}\leq
|f|_{\mathcal{B}^s_{p,q}(Q,\mathcal{P}_Q,\mathcal{A}_Q)}.$$
\begin{lstlisting}[style=lean]
def inducedRepresentationToAmbient
    (G : WeakGridSpace (α := α)) {k₀ : ℕ} (Q : LevelCell G k₀)
    {s : ℝ} {p u : ℝ≥0∞} [Fact (1 ≤ p)] (A : AtomFamily G s p u)
    {f : Lp ℂ p (inducedWeakGridSpace G Q).measure}
    (R : LpGridRepresentation (inducedAtomFamily G Q A) f) :
    LpGridRepresentation A (inducedLpToAmbient G Q f) := ...
\end{lstlisting}

\section{Examples of classes of atoms: Souza's atoms}\label{secatom}

Let $Q\in \mathcal{P}$. A {\bf $(s,p)$-Souza's atom} supported on $Q$ is a
function $a\colon I \rightarrow \mathbb{C}$ such that $a(x)=0$ for every
$x \not\in Q$ and $a$ is constant on $Q$, with
$$|a|_\infty\leq |Q|^{s-1/p}.$$
The set of Souza's atoms supported on $Q$ will be denoted by
$\mathcal{A}^{sz}_{s,p}(Q)$. The {\bf canonical Souza's atom} on $Q$ is the
Souza's atom with $a(x)= |Q|^{s-1/p}$ for every $x \in Q$. Souza's atoms are
$(s,p,\infty)$-type atoms and satisfy $A_1$--$A_5$.\footnote{Lean:
\lean{IsSouzaAtom}, \lean{canonicalSouzaAtom},
\lean{canonicalSouzaAtom\_isSouzaAtom} and the atom family
\lean{souzaAtomFamily} in \lean{GoodGrid/BesovSpace.lean}. The paper's
further examples --- H\"older atoms (Section~11B) and bounded variation atoms
(Section~11C) --- are not formalized and are omitted.}
\begin{lstlisting}[style=lean]
def IsSouzaAtom (G : GoodGridSpace (α := α)) (s : ℝ) (p : ℝ≥0∞)
    (Q : GoodGridCell G) (a : α → ℂ) : Prop :=
  (∀ x ∉ Q.cell, a x = 0) ∧
  ∃ c : ℂ, (∀ x ∈ Q.cell, a x = c) ∧
            ‖c‖ ≤ (G.grid.μ Q.cell).toReal ^ (s - (p.toReal)⁻¹)
\end{lstlisting}

\vspace{8mm}
\centerline{\Large\bf II. SPACES DEFINED BY SOUZA'S ATOMS}
\vspace{6mm}

\begin{center}
\fbox{\begin{minipage}{0.88\textwidth}
\centering
In Part II we suppose that $s>0$, \ $p\in [1,\infty)$ \ and \
$q\in [1,\infty]$.
\end{minipage}}
\end{center}
\vspace{4mm}

\section{Besov spaces in a measure space with a good grid}\label{secbesov}

We study the Besov-ish spaces
$\mathcal{B}^s_{p,q}=\mathcal{B}^s_{p,q}(\mathcal{P},\mathcal{A}^{sz}_{s,p})$
associated with a measure space with a good grid $(I,\mathcal{P},m)$. If
$p\in [1,\infty)$, $q\in [1,\infty]$ and $0<s<\frac{1}{p}$ we say that
$\mathcal{B}^s_{p,q}$ is a {\bf Besov space}. The general theory of Part~I
specializes as follows.

\begin{theorem}\label{souzamain}
Let $(I,\mathcal{P},m)$ be a measure space with a good grid, $s>0$,
$p \in [1,\infty)$, $q \in [1,\infty]$. Then
\begin{itemize}
\item[A.] $(\mathcal{B}^s_{p,q}, |\cdot|_{\mathcal{B}^s_{p,q}})$ is a complex
Banach space.\footnote{Lean: \lean{souzaBesovSpace\_costNorm\_completeSpace}
in \lean{GoodGrid/BesovSpace.lean}.}
\item[B.] Closed balls of $\mathcal{B}^s_{p,q}$ are compact subsets of $L^p$,
and also of $L^1$.\footnote{Lean: \lean{souza\_closedCostBallInLp\_isCompact}
and \lean{souza\_closedCostBallInL1\_isCompact} in
\lean{GoodGrid/BesovSpace.lean}. This is the Souza-atom instance of
Corollary~\ref{compa1} together with the embedding of
Proposition~\ref{lp}.}
\item[C.] If $0<s<1/p$ then $\mathcal{B}^s_{p,q}$ is dense in $L^p$, and
dense in $L^1$.\footnote{Lean: \lean{souzaBesovSpace\_dense} and
\lean{souzaBesovSpace\_dense\_inL1} in \lean{GoodGrid/BesovSpace.lean}.
The density statements do not appear explicitly as theorems in the paper;
they are recorded here because they are formalized.}
\end{itemize}
\end{theorem}
\begin{lstlisting}[style=lean,basicstyle=\ttfamily\scriptsize]
theorem souzaBesovSpace_costNorm_completeSpace
    (G : GoodGridSpace (α := α))
    (s : ℝ) (p q : ℝ≥0∞)
    (hs : 0 < s) (hp : 1 ≤ p) (hp_top : p ≠ ∞)
    [Fact (1 ≤ p)] [Fact (1 ≤ q)] :
    @CompleteSpace
      (SouzaBesovSpace G s p q hs hp hp_top)
      (WeakGridSpace.BesovishSpace.costNormedAddCommGroup
        (A := souzaAtomFamily G s p hs hp hp_top) (q := q)
        hp_top (souza_hCco G s p q hs hp hp_top)).toMetricSpace.toPseudoMetricSpace.toUniformSpace := ...

theorem souza_closedCostBallInLp_isCompact
    (G : GoodGridSpace (α := α))
    (s : ℝ) (p q : ℝ≥0∞)
    (hs : 0 < s) (hp : 1 ≤ p) (hp_top : p ≠ ∞)
    [Fact (1 ≤ p)] [Fact (1 ≤ q)]
    {C : ℝ} (hC : 0 ≤ C) :
    IsCompact
      {f : MeasureTheory.Lp ℂ p G.toWeakGridSpace.measure |
        ∃ hf : f ∈ SouzaBesovSpace G s p q hs hp hp_top,
          WeakGridSpace.BesovishSpace.Norm_Costpq
            (souzaAtomFamily G s p hs hp hp_top) q ⟨f, hf⟩ ≤ C} := ...

theorem souzaBesovSpace_dense
    (G : GoodGridSpace (α := α))
    (s : ℝ) (p q : ℝ≥0∞)
    (hs : 0 < s) (hp : 1 ≤ p) (hp_top : p ≠ ∞)
    [Fact (1 ≤ p)] [Fact (1 ≤ q)] :
    Dense (SouzaBesovSpace G s p q hs hp hp_top :
      Set (MeasureTheory.Lp ℂ p G.toWeakGridSpace.measure)) := ...
\end{lstlisting}

\section{Positive cone}\label{positivesec}

We say that $f$ is $\mathcal{B}^s_{p,q}$-positive if there is a
$\mathcal{B}^s_{p,q}$-representation
$$f= \sum_k \sum_{P\in \mathcal{P}^k} c_Pa_P$$
where $c_P\geq 0$ and $a_P$ is the canonical $(s,p)$-Souza's atom supported on
$P$. The set of all $\mathcal{B}^s_{p,q}$-positive functions is a convex cone
in $\mathcal{B}^s_{p,q}$, denoted $\mathcal{B}^{s+}_{p,q}$, with the
``norm''
$$|f|_{\mathcal{B}^{s+}_{p,q}}=\inf \Big( \sum_k \big( \sum_{P\in \mathcal{P}^k} c_P^p \big)^{q/p} \Big)^{1/q},$$
the infimum running over all positive representations of
$f$.\footnote{Lean: \lean{SouzaPositiveRepresentation} (canonical atoms,
nonnegative coefficients) and the cone \lean{souzaPositiveCone} in
\lean{GoodGrid/PositiveCone.lean}. The formalization also uses the wider
notion of a \emph{cone-positive} representation
(\lean{SouzaConePositiveRepresentation}): real coefficients $\geq 0$ and
Souza's atoms which are a.e.\ nonnegative real-valued, not necessarily
canonical.}
\begin{lstlisting}[style=lean]
def SouzaPositiveRepresentation
    (G : GoodGridSpace (α := α)) (s : ℝ) (p : ℝ≥0∞)
    (hs : 0 < s) (hp : 1 ≤ p) (hp_top : p ≠ ∞)
    [Fact (1 ≤ p)]
    {f : Lp ℂ p G.toWeakGridSpace.measure}
    (R : WeakGridSpace.LpGridRepresentation
      (souzaAtomFamily G s p hs hp hp_top) f) : Prop :=
  ∀ k, SouzaPositiveLevelBlock G s p hs hp hp_top (R.block k)

def SouzaConePositiveRepresentation
    (G : GoodGridSpace (α := α)) (s : ℝ) (p : ℝ≥0∞)
    (hs : 0 < s) (hp : 1 ≤ p) (hp_top : p ≠ ∞)
    [Fact (1 ≤ p)]
    {f : Lp ℂ p G.toWeakGridSpace.measure}
    (R : WeakGridSpace.LpGridRepresentation
      (souzaAtomFamily G s p hs hp hp_top) f) : Prop :=
  ∀ k, SouzaConePositiveLevelBlock G s p hs hp hp_top (R.block k)
\end{lstlisting}

\begin{proposition}\label{posdec}
Let $0 < s < 1/p$.
\begin{itemize}
\item[A.] If $f\in \mathcal{B}^s_{p,q}$ is (a.e.) real valued then one can
find $f_+,f_-\in \mathcal{B}^{s+}_{p,q}$ such that
$$f=f_+-f_-,$$
with the positive costs of $f_+$ and $f_-$ controlled in terms of
$|f|_{\mathcal{B}^{s}_{p,q}}$; an analogous four-term decomposition holds
for complex valued $f$.\footnote{Lean: \lean{exists\_souzaPositive\_decomposition\_of\_aeRealValued}
and \lean{exists\_souzaPositive\_decomposition\_of\_aeComplexValued} in
\lean{GoodGrid/PositiveCone.lean}. In the paper this appears as an unnumbered
observation in Section~13; the complex case is not stated in the paper.}
\item[B.] The positive cone $\mathcal{B}^{s+}_{p,q}$ is dense in the cone of
nonnegative functions of $L^p$.\footnote{Lean:
\lean{souzaPositiveCone\_dense\_in\_LpNonnegativeCone} in
\lean{GoodGrid/PositiveCone.lean}. Not stated in the paper.}
\end{itemize}
\end{proposition}
\begin{lstlisting}[style=lean,basicstyle=\ttfamily\scriptsize]
theorem exists_souzaPositive_decomposition_of_aeRealValued
    (G : GoodGridSpace (α := α)) (s : ℝ) (p q : ℝ≥0∞)
    (hs : 0 < s) (hp : 1 ≤ p) (hp_top : p ≠ ∞)
    [Fact (1 ≤ p)] [Fact (1 ≤ q)] :
    ∃ C : ℝ,
      0 ≤ C ∧
        ∀ f : WeakGridSpace.BesovishSpace
            (souzaAtomFamily G s p hs hp hp_top) q,
          SouzaBesovAEEqRealValued G s p q hs hp hp_top f →
            ∃ u v : WeakGridSpace.BesovishSpace
                (souzaAtomFamily G s p hs hp hp_top) q,
              SouzaPositiveElement G s p q hs hp hp_top u ∧
                SouzaPositiveElement G s p q hs hp hp_top v ∧
                f = u - v ∧
                souzaPositiveNorm G s p q hs hp hp_top u ≤
                  ENNReal.ofReal C *
                    ENNReal.ofReal (WeakGridSpace.BesovishSpace.Norm_Costpq
                      (souzaAtomFamily G s p hs hp hp_top) q f) ∧
                souzaPositiveNorm G s p q hs hp hp_top v ≤
                  ENNReal.ofReal C *
                    ENNReal.ofReal (WeakGridSpace.BesovishSpace.Norm_Costpq
                      (souzaAtomFamily G s p hs hp hp_top) q f) := ...

theorem souzaPositiveCone_dense_in_LpNonnegativeCone
    (G : GoodGridSpace (α := α)) (s : ℝ) (β q : ℝ≥0∞)
    (hs : 0 < s) (hβ : 1 ≤ β) (hβ_top : β ≠ ∞)
    [Fact (1 ≤ β)] [Fact (1 ≤ q)] :
    LpNonnegativeCone G β ⊆
      closure (SouzaPositiveConeInLbeta G s β q hs hβ hβ_top) := ...
\end{lstlisting}

\begin{proposition}[Support of positive functions]\label{possupp}
If $f \in \mathcal{B}^{s+}_{p,q}$ then its support
$$\operatorname{supp} f = \{ x\in I\colon \ f(x)\neq 0\}$$
is, up to a set of zero measure, a countable union of elements of
$\mathcal{P}$. Moreover, in any positive representation of $f$, the
coefficient of a cell $Q$ vanishes unless $Q$ is contained (mod $m$) in this
union.\footnote{Lean:
\lean{support\_ae\_countable\_iUnion\_goodGridCells\_of\_souzaPositiveFunction}
and the support theorem
\lean{souzaPositiveRepresentation\_coeff\_eq\_zero\_of\_not\_subset\_cell} in
\lean{GoodGrid/PositiveCone.lean} (paper Proposition~13.1; the second
sentence is a sharper statement obtained in the formalization).}
\end{proposition}
\begin{lstlisting}[style=lean]
theorem support_ae_countable_iUnion_goodGridCells_of_souzaPositiveFunction
    (G : GoodGridSpace (α := α)) (s : ℝ) (p q : ℝ≥0∞)
    (hs : 0 < s) (hp : 1 ≤ p) (hp_top : p ≠ ∞)
    [Fact (1 ≤ p)] [Fact (1 ≤ q)]
    {f : α → ℂ}
    (hf : SouzaPositiveFunction G s p q hs hp hp_top f) :
    IsAECountableUnionOfGoodGridCells G {x | f x ≠ 0} := ...

theorem souzaPositiveRepresentation_coeff_eq_zero_of_not_subset_cell
    (G : GoodGridSpace (α := α)) (s : ℝ) (p q : ℝ≥0∞)
    (hs : 0 < s) (hp : 1 ≤ p) (hp_top : p ≠ ∞)
    [Fact (1 ≤ p)] [Fact (1 ≤ q)]
    (Qc : GoodGridCell G)
    {g : Lp ℂ p G.toWeakGridSpace.measure}
    (R : WeakGridSpace.LpGridRepresentation (souzaAtomFamily G s p hs hp hp_top) g)
    (hR : SouzaPositiveRepresentation G s p hs hp hp_top R)
    (hsupp : ∀ᵐ x ∂ G.toWeakGridSpace.measure, x ∉ Qc.cell → (g : α → ℂ) x = 0)
    {n : ℕ} (P : WeakGridSpace.LevelCell G.toWeakGridSpace n)
    (hP : ¬ P.1 ⊆ Qc.cell) :
    (R.block n).coeff P = 0 := ...
\end{lstlisting}

\section{Unbalanced Haar wavelets}\label{haar}

Let $\mathcal{P}=(\mathcal{P}^k)_k$ be a good grid. Following Girardi and
Sweldens, one constructs, from the children structure of the grid (Type I
tree with the logarithmic subtrees construction), a family of {\bf unbalanced
Haar functions} $\{\phi_S\}_{S\in \hat{\mathcal{H}}}$, where
$\hat{\mathcal{H}}=\{I\}\cup \mathcal{H}$, $\phi_I = 1_I/|I|^{1/2}$, each
$\phi_S$ with $S \in \mathcal{H}_Q$ is supported on $Q$, has zero mean,
$L^2$-norm one, and satisfies
$$\frac{C_5(\lambda_1,\lambda_2)}{|Q|^{1/2}} \leq |\phi_S(x)| \leq \frac{C_6(\lambda_1,\lambda_2)}{|Q|^{1/2}}$$
on its support. The family $\{\phi_S\}_{S\in \hat{\mathcal{H}}}$ is an
unconditional basis of $L^\beta$ for every $1<\beta<
\infty$.\footnote{The construction and the unconditional basis property are
provided by the external Lean repositories \lean{UnbalancedHaarWavelet},
\lean{UnconditionalSchauderBasis}, \lean{LaminarFamiliesMaximalBinaryTrees}
and \lean{Burkholder} (SmaniaD), on which this repository depends. The local
API ($L^2$-normalized Haar functions \lean{L2normalizedHaar} and Haar
coefficients \lean{Coeff}) is in
\lean{GoodGrid/AlternativeRepresentationsAndNorms/HaarRepresentationNorm.lean}.}
\begin{lstlisting}[style=lean]
def L2normalizedHaar (G : GoodGridSpace (α := α)) [DecidableEq (Set α)]
    (F : UnbalancedHaarWavelet.FullHaarSystem (G := GridOf G))
    (i : F.Index) (x : α) : ℂ :=
  ((l2NormalizationFactor G F i : ℝ) : ℂ) *
    (UnbalancedHaarWavelet.FullHaarSystem.function (GridOf G) F i x : ℂ)

def Coeff (G : GoodGridSpace (α := α)) [DecidableEq (Set α)]
    (F : UnbalancedHaarWavelet.FullHaarSystem (G := GridOf G))
    (f : α → ℂ) (_hf : Integrable f G.grid.μ) (i : F.Index) : ℂ :=
  ∫ x, f x * L2normalizedHaar G F i x ∂G.grid.μ
\end{lstlisting}

\section{Alternative characterizations I: Messing with norms}\label{alt1}

For $f\in L^\beta$, $\beta>1$, write $d_S^f = \int f \phi_S\,dm$ for the Haar
coefficients of $f$, and define
$$ N_{haar}(f)=|I|^{1/p-s-1/2} |d_I^f|+ \Big( \sum_{k} \big( \sum_{Q\in \mathcal{P}^{k}} |Q|^{1-sp-\frac{p}{2}} \sum_{S\in \mathcal{H}_Q} |d_S^f|^p\big)^{q/p} \Big)^{1/q}.$$
From the Haar coefficients one builds, exactly as in the paper, the
{\bf standard atomic representation}
$$f = \sum_i \sum_{Q\in \mathcal{P}^{i}} k_Q^f\, a_Q$$ ($a_Q$ canonical Souza's
atoms) and its norm
$$N_{st}(f)=|k_I^f|+ \Big( \sum_{k\geq 1} \big( \sum_{Q\in \mathcal{P}^{k}} |k_Q^f|^p \big)^{q/p}\Big)^{1/q}.$$
Finally, with
$$osc_p(f,Q)= \inf_{c \in \mathbb{C}} \Big(\int_Q |f-c|^p \, dm\Big)^{1/p},$$
define
$$N_{osc}(f)=|I|^{-s}|f|_p+\Big( \sum_k \big( \sum_{Q\in \mathcal{P}^k} |Q|^{-sp} osc_p(f,Q)^p \big)^{q/p} \Big)^{1/q}.$$

\begin{theorem}[Equivalence of norms]\label{alte}
Suppose $0<s<1/p$, $p\in [1,\infty)$ and $q\in [1,\infty]$. Each one of the
quantities $|f|_{\mathcal{B}^s_{p,q}}$, $N_{st}(f)$, $N_{haar}(f)$,
$N_{osc}(f)$ is finite if and only if $f\in \mathcal{B}^s_{p,q}$, and these
norms are equivalent on $\mathcal{B}^s_{p,q}$. More precisely there are
constants, depending only on the grid geometry and on $s,p,q$, such that
\begin{align*}
N_{st}(f) &\leq C\, |f|_{\mathcal{B}^s_{p,q}}, &
N_{st}(f) &\leq C\, N_{haar}(f),\\
N_{haar}(f) &\leq C\, N_{osc}(f), &
N_{osc}(f) &\leq C\, |f|_{\mathcal{B}^s_{p,q}},
\end{align*}
and conversely: if $N_{st}(f)<\infty$ then $f \in \mathcal{B}^s_{p,q}$ and
the standard atomic representation is a
$\mathcal{B}^s_{p,q}$-representation of $f$ (so
$|f|_{\mathcal{B}^s_{p,q}} \leq N_{st}(f)$); if $N_{haar}(f) < \infty$ and
$f \in L^1$ then $f \in L^p$ and the Haar series of $f$ converges in
$L^p$.\footnote{Lean: the inequality cycle
\lean{exists\_standardRepresentationNorm\_le\_const\_mul\_souzaBesovNorm},
\lean{exists\_standardRepresentationNorm\_le\_const\_mul\_haarL2RepresentationNorm},
\lean{exists\_haarL2RepresentationNorm\_le\_const\_mul\_meanOscillationNorm},
\lean{exists\_meanOscillationNorm\_le\_const\_mul\_souzaBesovNorm},
\lean{finite\_standardRepresentationNorm\_implies\_}\allowbreak\lean{memBesov\_and\_standardRepresentation},
\lean{finite\_haarL2RepresentationNorm\_implies\_memLp\_and\_hasSum} and
\lean{exists\_haarRepresentationNorms\_equivalent}, in the files of
\lean{GoodGrid/AlternativeRepresentationsAndNorms/}. Compared with paper
Theorem~15.1: the formalization proves the equivalences as a cycle of
separate inequalities with existential constants (the explicit constants
$\theta_1$, $\theta_2$ of the paper are not tracked), it works with the
$L^2$-normalized Haar coefficients, and there is no single wrap-up statement
packaging all four norms at once.}
\end{theorem}
\begin{lstlisting}[style=lean,basicstyle=\ttfamily\scriptsize]
theorem exists_standardRepresentationNorm_le_const_mul_haarL2RepresentationNorm
    (G : GoodGridSpace (α := α)) [DecidableEq (Set α)]
    (F : UnbalancedHaarWavelet.FullHaarSystem (G := HaarRepresentation.GridOf G))
    [DecidableEq F.Index]
    (s : ℝ) (hs : 0 < s)
    (p : ℝ≥0∞) [Fact (1 ≤ p)] (hp_top : p < ∞)
    (q : ℝ≥0∞) [Fact (1 ≤ q)] :
    ∃ C : ℝ≥0∞, C ≠ ∞ ∧
      ∀ (f : α → ℂ) (hfint : Integrable f G.grid.μ),
        HaarRepresentation.haarL2RepresentationNorm G F s p q f hfint ≠ ∞ →
          standardRepresentationNorm G F s hs p hp_top q f hfint ≠ ∞ ∧
            standardRepresentationNorm G F s hs p hp_top q f hfint ≤
              C * HaarRepresentation.haarL2RepresentationNorm G F s p q f hfint := ...

theorem exists_haarL2RepresentationNorm_le_const_mul_meanOscillationNorm
    (G : GoodGridSpace (α := α)) [DecidableEq (Set α)]
    (F : UnbalancedHaarWavelet.FullHaarSystem (G := GridOf G))
    (s : ℝ) (p : ℝ≥0∞) [Fact (1 ≤ p)] (hp_top : p < ∞)
    (q : ℝ≥0∞) [Fact (1 ≤ q)] :
    ∃ C : ℝ≥0∞, C ≠ ∞ ∧
      ∀ (f : α → ℂ) (hf : Integrable f G.grid.μ), MemLp f p G.grid.μ →
        MeanOscillation.meanOscillationNorm G s p q f ≠ ∞ →
          haarL2RepresentationNorm G F s p q f hf ≤
            C * MeanOscillation.meanOscillationNorm G s p q f := ...

theorem exists_meanOscillationNorm_le_const_mul_souzaBesovNorm
    (G : GoodGridSpace (α := α)) [DecidableEq (Set α)]
    (s : ℝ) (hs : 0 < s)
    (p : ℝ≥0∞) [Fact (1 ≤ p)] (hp_top : p < ∞)
    (q : ℝ≥0∞) [Fact (1 ≤ q)] :
    ∃ C : ℝ≥0∞, C ≠ ∞ ∧
      ∀ (g : SouzaBesovSpace G s p q hs Fact.out (ne_of_lt hp_top))
        (f : α → ℂ),
          f =ᵐ[G.grid.μ] ((g : Lp ℂ p G.toWeakGridSpace.measure) : α → ℂ) →
            meanOscillationNorm G s p q f ≤
              C * ENNReal.ofReal
                (WeakGridSpace.BesovishSpace.Norm_Costpq
                  (souzaAtomFamily G s p hs Fact.out (ne_of_lt hp_top)) q g) := ...

theorem finite_standardRepresentationNorm_implies_memBesov_and_standardRepresentation
    (G : GoodGridSpace (α := α)) [DecidableEq (Set α)]
    (F : UnbalancedHaarWavelet.FullHaarSystem (G := HaarRepresentation.GridOf G))
    [DecidableEq F.Index]
    (s : ℝ) (hs : 0 < s)
    (p : ℝ≥0∞) [Fact (1 ≤ p)] (hp_top : p < ∞)
    (q : ℝ≥0∞) [Fact (1 ≤ q)]
    (f : α → ℂ) (hf : Integrable f G.grid.μ)
    (hN : standardRepresentationNorm G F s hs p hp_top q f hf ≠ ∞) :
    ∃ hfLp : MemLp f p G.grid.μ,
      ∃ g : SouzaBesovSpace G s p q hs Fact.out (ne_of_lt hp_top),
        ∃ R :
          WeakGridSpace.LpGridRepresentation
            (souzaAtomFamily G s p hs Fact.out (ne_of_lt hp_top))
            (g : Lp ℂ p G.toWeakGridSpace.measure),
          (g : Lp ℂ p G.toWeakGridSpace.measure) = hfLp.toLp f ∧
            R.block = canonicalStandardBlockSeq G F s hs p hp_top f hf ∧
              WeakGridSpace.LpGridRepresentation.FinitePQCost (q := q) R ∧
                WeakGridSpace.LpGridRepresentation.pqCostENNReal (q := q) R =
                  standardRepresentationNorm G F s hs p hp_top q f hf ∧
                  WeakGridSpace.LpGridRepresentation.pqCost (q := q) R ≤
                    (standardRepresentationNorm G F s hs p hp_top q f hf).toReal ∧
                    WeakGridSpace.BesovishSpace.Norm_Costpq
                        (souzaAtomFamily G s p hs Fact.out (ne_of_lt hp_top)) q g ≤
                      (standardRepresentationNorm G F s hs p hp_top q f hf).toReal := ...
\end{lstlisting}

\begin{corollary}[Reading the standard coefficients]\label{fou}
For each $P\in \mathcal{P}$ the coefficient $k_P^f$ of the standard atomic
representation of $f$ is a (finite) linear combination of the Haar
coefficients $d^f_S$, $S \in \mathcal H_Q$, of $f$,\footnotemark\ and the standard
representation of any $f \in \mathcal{B}^s_{p,q}$ converges to $f$ in
$L^p$ with
$$\Big(\sum_k \big( \sum_{P\in \mathcal{P}^k} |k_P^f|^p\big)^{q/p}\Big)^{1/q}\leq C\, |f|_{\mathcal{B}^s_{p,q}}.$$
\footnotetext{Lean: \lean{Coeff} and
\lean{hasSum\_coeff\_smul\_l2normalizedHaar\_toLp} in
\lean{GoodGrid/AlternativeRepresentationsAndNorms/HaarRepresentationNorm.lean},
together with the theorems of the previous footnote. Paper Corollary~15.2
states moreover that $f \mapsto k^f_P$ extends to a bounded linear functional
on $L^1$; this extension is not formalized.}
\end{corollary}
\begin{lstlisting}[style=lean]
theorem hasSum_coeff_smul_l2normalizedHaar_toLp
    (G : GoodGridSpace (α := α)) [DecidableEq (Set α)]
    (F : UnbalancedHaarWavelet.FullHaarSystem (G := GridOf G))
    [DecidableEq F.Index]
    (β : ℝ≥0∞) (hβ_one : 1 < β) (hβ_top : β < ∞)
    (f : α → ℂ) (hf : MemLp f β G.grid.μ) :
    letI : Fact (1 ≤ β) := ⟨le_of_lt hβ_one⟩
    HasSum
      (fun i : F.Index =>
        Coeff G F f
            (by
              letI : IsFiniteMeasure G.grid.μ := (GridOf G).isFinite
              exact hf.integrable (le_of_lt hβ_one))
            i •
          (l2normalizedHaar_memLp G F β i).toLp (L2normalizedHaar G F i))
      (hf.toLp f) := ...
\end{lstlisting}

\section{Alternative characterizations II: Messing with atoms}\label{alt2}

Let $s < \beta$ and $\tilde{q}\in [1,\infty]$. A
{\bf $(s,\beta,p,\tilde{q})$-Besov atom} on $Q$ is a function
$a \in \mathcal{B}^\beta_{p,\tilde{q}}(Q,\mathcal{P}_Q,\mathcal{A}^{sz}_{\beta,p})$
such that $a(x)=0$ for $x \not\in Q$ and
$$|a|_{\mathcal{B}^\beta_{p,\tilde{q}}(Q, \mathcal{P}_Q, \mathcal{A}^{sz}_{\beta,p})}\leq \frac{1}{C_7} |Q|^{s-\beta},$$
where $C_7 = C_7(C_1,\lambda_2,\beta,\tilde q,p) \geq 1$ is the normalizing
constant of the paper. The family of $(s,\beta,p,\tilde{q})$-Besov atoms
supported on $Q$ is denoted $\mathcal{A}^{bs}_{s,\beta,p,\tilde{q}}(Q)$, and
$\mathcal{A}^{sz}_{s,p}(Q)\subset \mathcal{A}^{bs}_{s,\beta,p,\tilde{q}}(Q)$;
Besov atoms are atoms of type $(s,p,1)$.

\begin{proposition}[Souza's atoms and Besov's atoms]\label{besova}
Let $\mathcal{P}$ be a good grid and $0<s<1/p$. Let $\mathcal{A}$ be a class
of $(s,p,u)$-atoms, with $u\geq 1$, such that for some $s< \beta$,
$\tilde{q}\in [1,\infty]$, and $C_8, C_9 \geq 0$ we have, for every
$Q\in \mathcal{P}$,
$$\frac{1}{C_8} \mathcal{A}^{sz}_{s,p}(Q) \subset \mathcal{A}(Q)\subset C_9\, \mathcal{A}^{bs}_{s,\beta,p,\tilde{q}}(Q).$$
Then\footnotemark
$$\mathcal{B}^s_{p,q}(\mathcal{P},\mathcal{A}^{sz}_{s,p})=\mathcal{B}^s_{p,q}(\mathcal{P},\mathcal{A})=\mathcal{B}^s_{p,q}(\mathcal{P},\mathcal{A}^{bs}_{s,\beta,p,\tilde{q}}),$$
with
$$|f|_{\mathcal{B}^s_{p,q}(\mathcal{A})} \leq C_8\, |f|_{\mathcal{B}^s_{p,q}(\mathcal{A}^{sz}_{s,p})}
\quad\text{and}\quad
|f|_{\mathcal{B}^s_{p,q}(\mathcal{A}^{sz}_{s,p})} \leq C\, |f|_{\mathcal{B}^s_{p,q}(\mathcal{A})}.$$
\footnotetext{Lean: \lean{atoms\_between\_souza\_atoms\_and\_besov\_atoms} in
\lean{GoodGrid/BesovAtoms.lean}, supported by the decay claims
\lean{besovAtom\_to\_souza\_representation\_decay} and
\lean{besovAtom\_to\_induced\_souzaS\_representation\_decay\_claimC} (the
Lean counterparts of the Claim in the paper's proof, with the geometric decay
rate $\lambda_2^{(k-j)(\beta-s)}$). Paper Propositions~16.2 (H\"older atoms)
and 16.3 (bounded variation atoms) are not formalized and are omitted.}
\end{proposition}
\begin{lstlisting}[style=lean,basicstyle=\ttfamily\scriptsize]
theorem atoms_between_souza_atoms_and_besov_atoms
    (G : GoodGridSpace (α := α)) (s β : ℝ) (p u q qtilde : ℝ≥0∞)
    [Fact (1 ≤ p)] [Fact (1 ≤ u)] [Fact (1 ≤ q)] [Fact (1 ≤ qtilde)]
    (hs : 0 < s) (hβ : 0 < β) (hβs : s < β) (hp_top : p ≠ ∞)
    (A : WeakGridSpace.AtomFamily G.toWeakGridSpace s p u)
    (C1 C2 : ℝ)
    (hSandwich :
      SouzaBesovSandwich G s β p u qtilde hs hβ (inferInstance : Fact (1 ≤ p))
        hp_top A C1 C2) :
    (WeakGridSpace.BesovishSpace (souzaAtomFamily G s p hs (Fact.out : 1 ≤ p) hp_top) q =
        WeakGridSpace.BesovishSpace A q) ∧
      (WeakGridSpace.BesovishSpace A q =
        WeakGridSpace.BesovishSpace
          (besovAtomFamily G s β p qtilde hs hβ inferInstance hp_top) q) ∧
      (∀ f : WeakGridSpace.BesovishSpace
          (souzaAtomFamily G s p hs (Fact.out : 1 ≤ p) hp_top) q,
        ∃ hfA : WeakGridSpace.MemBesovishCoeffCost A q
            (f : Lp ℂ p G.toWeakGridSpace.measure),
          WeakGridSpace.BesovishSpace.Norm_Costpq A q
              (⟨(f : Lp ℂ p G.toWeakGridSpace.measure), hfA⟩ :
                WeakGridSpace.BesovishSpace A q)
            ≤ C1 *
              WeakGridSpace.BesovishSpace.Norm_Costpq
                (souzaAtomFamily G s p hs (Fact.out : 1 ≤ p) hp_top) q f) ∧
      (∀ f : WeakGridSpace.BesovishSpace A q,
        ∃ hfS : WeakGridSpace.MemBesovishCoeffCost
            (souzaAtomFamily G s p hs (Fact.out : 1 ≤ p) hp_top) q
            (f : Lp ℂ p G.toWeakGridSpace.measure),
          WeakGridSpace.BesovishSpace.Norm_Costpq
              (souzaAtomFamily G s p hs (Fact.out : 1 ≤ p) hp_top) q
              (⟨(f : Lp ℂ p G.toWeakGridSpace.measure), hfS⟩ :
                WeakGridSpace.BesovishSpace
                  (souzaAtomFamily G s p hs (Fact.out : 1 ≤ p) hp_top) q)
            ≤ (C2 / (1 - G.grid.lambda2 ^ (β - s))) *
              WeakGridSpace.BesovishSpace.Norm_Costpq A q f) := ...
\end{lstlisting}

\vspace{8mm}
\centerline{\Large\bf III. APPLICATIONS}
\vspace{6mm}

\begin{center}
\fbox{\begin{minipage}{0.88\textwidth}
\centering
In Part III we suppose that $0<s<1/p$, \ $p\in [1,\infty)$ \ and \
$q\in [1,\infty]$.
\end{minipage}}
\end{center}
\vspace{4mm}

\section{Pointwise multipliers acting on $\mathcal{B}^s_{p,q}$}\label{multsec}

Let $g\colon I \rightarrow \mathbb{C}$ be a measurable function. We say that
$g$ is a {\bf pointwise multiplier} acting on $\mathcal{B}^s_{p,q}$ if the
transformation $G(f)=gf$ defines a bounded operator in $\mathcal{B}^s_{p,q}$.
We denote the set of pointwise multipliers by $M(\mathcal{B}^s_{p,q})$, with
the norm
$$|g|_{M(\mathcal{B}^s_{p,q})}=\sup\{|gf|_{\mathcal{B}^s_{p,q}}\colon \ |f|_{\mathcal{B}^s_{p,q}}\leq 1 \}.$$
Define
$$\mathcal{B}^s_{p,q,selfs}=\Big\{ g\in \mathcal{B}^s_{p,q}\colon \ \sup_{a_Q \in \mathcal{A}^{sz}_{s,p}} |ga_Q|_{\mathcal{B}^s_{p,q}} < \infty\Big\},$$
with the norm
$$|g|_{\mathcal{B}^s_{p,q,selfs}}=\sup_{a_Q \in \mathcal{A}^{sz}_{s,p}} |ga_Q|_{\mathcal{B}^s_{p,q}},$$
and, for $t \in \mathbb{N}$, the tail seminorms\footnote{Lean: \lean{SouzaPointwiseMultiplier},
\lean{SouzaPointwiseSelfsClass} and the tail class
\lean{SouzaPointwiseSelfsTailClass} with their seminorms, in
\lean{GoodGrid/Multipliers/Definition.lean}.}
$$|g|_{\mathcal{B}^{s,t}_{p,q,selfs}}=\sup_{\substack{ a_Q \in \mathcal{A}^{sz}_{s,p}\\ Q\in \mathcal{P}^j,\ j\geq t}} |ga_Q|_{\mathcal{B}^s_{p,q}}.$$
\begin{lstlisting}[style=lean]
abbrev SouzaPointwiseMultiplier
    (G : GoodGridSpace (α := α)) (s : ℝ) (p q : ℝ≥0∞)
    (hs : 0 < s) (hp : 1 ≤ p) (hp_top : p ≠ ∞)
    [Fact (1 ≤ p)] [Fact (1 ≤ q)]
    (m : α → ℂ) : Prop :=
  WeakGridSpace.IsPointwiseMultiplier
    (A := souzaAtomFamily G s p hs hp hp_top) q m

abbrev SouzaPointwiseSelfsClass
    (G : GoodGridSpace (α := α)) (s : ℝ) (p q : ℝ≥0∞)
    (hs : 0 < s) (hp : 1 ≤ p) (hp_top : p ≠ ∞)
    [Fact (1 ≤ p)] [Fact (1 ≤ q)]
    (m : α → ℂ) : Prop :=
  WeakGridSpace.PointwiseSelfsClass
    (A := souzaAtomFamily G s p hs hp hp_top) q m
\end{lstlisting}

\begin{proposition}\label{m11}
We have that\footnotemark
$$M(\mathcal{B}^s_{1,1})=\mathcal{B}^s_{1,1,selfs}.$$
\footnotetext{Lean:
\lean{souzaPointwiseMultiplier\_iff\_souzaPointwiseSelfsClass\_one\_one}
(the inclusion $\mathcal{B}^s_{1,1,selfs}\subset M(\mathcal{B}^s_{1,1})$ is
proved in \lean{GoodGrid/Multipliers/Besovs11.lean}; the converse, which
holds for all $p,q$, is
\lean{souzaPointwiseSelfsClass\_of\_souzaPointwiseMultiplier} in
\lean{GoodGrid/Multipliers/Besovspq.lean}). This is paper
Proposition~18.1.}
\end{proposition}
\begin{lstlisting}[style=lean]
theorem souzaPointwiseMultiplier_iff_souzaPointwiseSelfsClass_one_one
    (G : GoodGridSpace (α := α)) (s : ℝ)
    (hs : 0 < s) {m : α → ℂ} :
    SouzaPointwiseMultiplier G s (1 : ℝ≥0∞) (1 : ℝ≥0∞)
      hs le_rfl ENNReal.one_ne_top m ↔
    SouzaPointwiseSelfsClass G s (1 : ℝ≥0∞) (1 : ℝ≥0∞)
      hs le_rfl ENNReal.one_ne_top m := ...
\end{lstlisting}

\begin{lemma}[Restriction to a cell]\label{incint3}
Let $W\in \mathcal{P}$. The restriction map
$$r\colon \mathcal{B}^s_{p,q}(I, \mathcal{P}, \mathcal{A}^{sz}_{s,p})\rightarrow \mathcal{B}^s_{p,q}(W, \mathcal{P}_W, \mathcal{A}^{sz}_{s,p}),
\qquad r(f)=1_Wf,$$
is continuous: there is $C_{10} \geq 1$, independent of $W$, such that
\begin{itemize}
\item[A.] For every $f \in \mathcal{B}^s_{p,q}$ we have
$$|1_Wf|_{\mathcal{B}^s_{p,q}(W, \mathcal{P}_W, \mathcal{A}^{sz}_{s,p})}\leq C_{10}\, |f|_{\mathcal{B}^s_{p,q}}.$$
In particular $1_W\in M(\mathcal{B}^s_{p,q})$ with
$|1_W|_{M(\mathcal{B}^s_{p,q})} \le C_{10}$, uniformly in $W$.
\item[B.] For every $\mathcal{B}^{s+}_{p,q}$-representation
$f=\sum_{k}\sum_{Q} c_Qa_Q$ one can find a positive representation
$$1_Wf= \sum_{k}\sum_{Q} d_Q\,a_Q$$
on the induced grid with cost at most
$C_{10}$ times the cost of the original one, and such that $d_Q\neq 0$
implies $Q\subset \operatorname{supp} f$ (mod $m$).\footnote{Lean: \lean{souzaCellIndicatorRestrictsToInduced}, the bounds
\lean{souzaIndicatorRestrictionBound\_of\_ambientRestrictionTransmutation\_*},
\lean{souzaCellIndicatorPointwiseMultiplierBound\_uniform} and
\lean{souzaCellIndicatorPointwiseMultiplier} in
\lean{GoodGrid/Multipliers/Besovspq.lean} (paper Lemma~18.2). The paper's
strongly regular domains (Definition~18.7 and Proposition~18.8), of which the
present lemma is the grid-cell case, are not formalized in general; the
multiplier statement for $1_W$, $W \in \mathcal{P}$, with a bound uniform in
$W$, is the formalized special case.}
\end{itemize}
\end{lemma}
\begin{lstlisting}[style=lean]
theorem souzaCellIndicatorPointwiseMultiplier
    (G : GoodGridSpace (α := α)) (s : ℝ) (p q : ℝ≥0∞)
    (hs : 0 < s) (hp : 1 ≤ p) (hp_top : p ≠ ∞)
    [Fact (1 ≤ p)] [Fact (1 ≤ q)]
    (hs_lt_inv : s < (p.toReal)⁻¹)
    (W : GoodGridCell G) :
    SouzaPointwiseMultiplier G s p q hs hp hp_top
      (W.cell.indicator fun _ => (1 : ℂ)) := ...

theorem souzaCellIndicatorPointwiseMultiplierBound_uniform
    (G : GoodGridSpace (α := α)) (s : ℝ) (p q : ℝ≥0∞)
    (hs : 0 < s) (hp : 1 ≤ p) (hp_top : p ≠ ∞)
    [Fact (1 ≤ p)] [Fact (1 ≤ q)]
    (hs_lt_inv : s < (p.toReal)⁻¹)
    (W : GoodGridCell G) :
    SouzaPointwiseMultiplierBound G s p q hs hp hp_top
      (W.cell.indicator fun _ => (1 : ℂ))
      (2 * souzaAmbientRestrictionMultiplierConstant G s p + 1) := ...
\end{lstlisting}

\begin{proposition}[$selfs$ classes are bounded]\label{selfsbdd}
There is $C_{11}=C_{11}(\lambda_1,\lambda_2,C_1,s,p,q)$ such that for every
$t \in \mathbb{N}$ and every $g$ with
$|g|_{\mathcal{B}^{s,t}_{p,q,selfs}}<\infty$,
$$|g(x)| \leq C_{11}\, |g|_{\mathcal{B}^{s,t}_{p,q,selfs}}
\quad \text{for } m\text{-a.e. } x.$$
In particular $\mathcal{B}^s_{p,q,selfs} \subset L^\infty$ and this inclusion
is continuous.\footnote{Lean:
\lean{souzaPointwiseSelfsTailBound\_norm\_ae\_le},
\lean{souzaPointwiseSelfsTailNorm\_norm\_ae\_le} and
\lean{souzaPointwiseSelfsTailClass\_norm\_ae\_bounded} in
\lean{GoodGrid/Multipliers/MultipliersareBounded.lean}. This sharpens paper
Proposition~18.3 in two respects: the estimate is proved for the tail
seminorms $|\cdot|_{\mathcal{B}^{s,t}_{p,q,selfs}}$ with a constant
\emph{independent of $t$} (this uniformity is what the non-Archimedean
arguments below require), and the formalization also provides a
\emph{canonical} variant
(\lean{SouzaPointwiseCanonicalSelfsTailBound}) in which $g$ is only tested
against canonical Souza's atoms.}
\end{proposition}
\begin{lstlisting}[style=lean]
theorem souzaPointwiseSelfsTailNorm_norm_ae_le
    (G : GoodGridSpace (α := α))
    (s : ℝ) (p q : ℝ≥0∞)
    (hs : 0 < s) (hp : 1 ≤ p) (hp_top : p ≠ ∞)
    [Fact (1 ≤ p)] [Fact (1 ≤ q)]
    (hs_lt_inv : s < (p.toReal)⁻¹)
    {t : ℕ} {m : α → ℂ}
    (hm : SouzaPointwiseSelfsTailClass G s p q hs hp hp_top t m) :
    ∀ᵐ x ∂G.grid.μ,
      ‖m x‖ ≤
        souzaBesovLpLocalEmbeddingConstant G s p q *
          (2 * souzaAmbientRestrictionMultiplierConstant G s p + 1) *
            souzaPointwiseSelfsTailNorm G s p q hs hp hp_top t m := ...
\end{lstlisting}

\subsection{Non-Archimedean behaviour in $\mathcal{B}^\beta_{p,\tilde{q},selfs}$}

If $g_i \in M(\mathcal{B}^s_{p,q})$, the naive estimate is
$$\Big|\big(\sum_i g_i\big)f\Big|_{\mathcal{B}^s_{p,q}}\leq \Big(\sum_i
|g_i|_{M(\mathcal{B}^s_{p,q})}\Big)\,|f|_{\mathcal{B}^s_{p,q}}.$$
Remarkably, when the supports of the $g_i$ are separated by the grid, one
gets a far better estimate.

\begin{theorem}[Non-Archimedean property, finite version]\label{sepa}
Let $\beta > s$ and $\tilde q \in [1,\infty]$. There is
$C_{12}=C_{12}(\lambda_1,\lambda_2,C_1,s,\beta,p,q,\tilde q)$ with the
following property. Let $\Lambda$ be finite, $g_i$ with
$|g_i|_{\mathcal{B}^{\beta,t_i}_{p,\tilde{q},selfs}} < \infty$ for $i\in
\Lambda$ and $t_i \in \mathbb{N}$. Consider a function $f$ with a
$\mathcal{B}^s_{p,q}$-representation
$$ f=\sum_k \sum_{Q\in \mathcal{P}^k} c_Q a_Q$$
by canonical Souza's atoms satisfying
\begin{itemize}
\item[A.] We have
$$\sup_{\substack{ Q\in \mathcal{P} \\ c_Q\neq 0}} \
\sum_{i\colon\, Q\cap \operatorname{supp} g_i \neq \emptyset} |g_i|_{ \mathcal{B}^{\beta,t_i}_{p,\tilde{q},selfs}} \leq N.$$
\item[B.] If $Q\in \mathcal{P}^k$ satisfies $c_Q\neq 0$ and
$Q\cap \operatorname{supp} g_i \neq \emptyset$ then $k\geq t_i$.
\end{itemize}
Then $(\sum_{i\in\Lambda} g_i)f \in \mathcal{B}^s_{p,q}$ and we can find a
$\mathcal{B}^s_{p,q}$-representation\footnotemark
\begin{equation} \label{rre} \big(\sum_{i\in\Lambda} g_i\big)f= \sum_k \sum_{Q\in \mathcal{P}^k} d_Q a_Q\end{equation}
by canonical Souza's atoms such that
\begin{equation}\label{reep33} \Big(\sum_k \big( \sum_{Q\in \mathcal{P}^k} |d_Q|^p \big)^{q/p} \Big)^{1/q}
\leq C_{12}\, N\, \Big(\sum_k \big( \sum_{Q\in \mathcal{P}^k} |c_Q|^p \big)^{q/p} \Big)^{1/q}.
\end{equation}
\footnotetext{Lean: \lean{souzaNonArchimedeanPropertyLambdaFinite} in
\lean{GoodGrid/Multipliers/NonArchimedeanProperty.lean} (paper
Proposition~18.4 for finite $\Lambda$). The representations of $f$ and of
the product are by \emph{canonical} Souza's atoms
(\lean{SouzaCanonicalRepresentation}), which is no loss of generality for
Souza's atoms and is the precise form formalized.}
\end{theorem}
\begin{lstlisting}[style=lean,basicstyle=\ttfamily\scriptsize]
theorem souzaNonArchimedeanPropertyLambdaFinite
    (G : GoodGridSpace (α := α))
    (s β : ℝ) (p q qtilde : ℝ≥0∞)
    (hs : 0 < s) (hβ : 0 < β) (hβs : s < β)
    (hβ_lt_inv : β < (p.toReal)⁻¹)
    (hp : 1 ≤ p) (hp_top : p ≠ ∞)
    [Fact (1 ≤ p)] [Fact (1 ≤ q)] [Fact (1 ≤ qtilde)] :
    ∃ Cgen : ℝ,
      0 ≤ Cgen ∧
      ∀ (Λ : Finset ℕ) (t : ℕ → ℕ) (g : ℕ → α → ℂ) (N : ℝ)
        (f : α → ℂ)
        (x : WeakGridSpace.BesovishSpace
          (souzaAtomFamily G s p hs hp hp_top) q)
        (R : WeakGridSpace.LpGridRepresentation
          (souzaAtomFamily G s p hs hp hp_top)
          (x : Lp ℂ p G.toWeakGridSpace.measure)),
          0 ≤ N →
          WeakGridSpace.RepresentsFunction
            (G := G.toWeakGridSpace) (p := p) f
            (x : Lp ℂ p G.toWeakGridSpace.measure) →
          WeakGridSpace.LpGridRepresentation.FinitePQCost (q := q) R →
          (∀ i ∈ Λ,
            ∃ C : ℝ,
              SouzaPointwiseSelfsTailBound
                G β p qtilde hβ hp hp_top (t i) (g i) C) →
          (∀ k (Q : WeakGridSpace.LevelCell G.toWeakGridSpace k),
            (R.block k).coeff Q ≠ 0 →
              nonArchimedeanRelevantTailSelfsSum
                G β p qtilde hβ hp hp_top Λ t g Q ≤ N) →
          (∀ k (Q : WeakGridSpace.LevelCell G.toWeakGridSpace k) i,
            i ∈ Λ →
              (R.block k).coeff Q ≠ 0 →
                goodGridLevelCellMeetsSupport G Q (g i) →
                  t i ≤ k) →
          ∃ y : WeakGridSpace.BesovishSpace
              (souzaAtomFamily G s p hs hp hp_top) q,
            ∃ S : WeakGridSpace.LpGridRepresentation
                (souzaAtomFamily G s p hs hp hp_top)
                (y : Lp ℂ p G.toWeakGridSpace.measure),
              WeakGridSpace.RepresentsFunction
                (G := G.toWeakGridSpace) (p := p)
                (fun z => (∑ i ∈ Λ, g i z) * f z)
                (y : Lp ℂ p G.toWeakGridSpace.measure) ∧
              WeakGridSpace.LpGridRepresentation.pqCost (q := q) S ≤
                Cgen * N *
                  WeakGridSpace.LpGridRepresentation.pqCost (q := q) R := ...
\end{lstlisting}

\begin{theorem}[Non-Archimedean property, infinite version]\label{sepainf}
In the setting of Theorem~\ref{sepa} let $\Lambda \subset \mathbb{N}$ be
infinite, and replace hypothesis A.\ by the summability condition
$$\sum_{i \in \Lambda} T_i \leq N, \qquad
T_i := \sup_{\substack{ Q\in \mathcal{P},\ c_Q\neq 0 \\ Q\cap \operatorname{supp} g_i \neq \emptyset}}
|g_i|_{ \mathcal{B}^{\beta,t_i}_{p,\tilde{q},selfs}}\ ,$$
understood as a series with values in $[0,\infty]$. Then, for $m$-a.e.\ $x$
in $\{f \neq 0\}$, the series $\sum_{i\in\Lambda} g_i(x)$ converges
absolutely, with
$$\sum_{i\in\Lambda} |g_i(x)| \leq C_{12}\, N;$$
the function
$h = (\sum_{i\in\Lambda} g_i) f$ (defined a.e.\ as this pointwise limit
times $f$) belongs to $L^p$ and to $\mathcal{B}^s_{p,q}$, and it has a
canonical representation (\ref{rre}) satisfying
(\ref{reep33}).\footnote{Lean: \lean{souzaNonArchimedeanProperty} in
\lean{GoodGrid/Multipliers/NonArchimedeanProperty.lean}. Paper
Proposition~18.4 treats finite and infinite $\Lambda$ together, reducing to
the finite case; the formalization keeps them separate because the infinite
case has genuinely more content (pointwise absolute convergence, membership
of the limit in $L^p$, and identification of the limit, obtained through the
compactness of representations with uniform cost,
\lean{exists\_strongly\_convergent\_subseq\_of\_uniform\_pqCost}).}
\end{theorem}
\begin{lstlisting}[style=lean,basicstyle=\ttfamily\scriptsize]
theorem souzaNonArchimedeanProperty
    (G : GoodGridSpace (α := α))
    (s β : ℝ) (p q qtilde : ℝ≥0∞)
    (hs : 0 < s) (hβ : 0 < β) (hβs : s < β)
    (hβ_lt_inv : β < (p.toReal)⁻¹)
    (hp : 1 ≤ p) (hp_top : p ≠ ∞)
    [Fact (1 ≤ p)] [Fact (1 ≤ q)] [Fact (1 ≤ qtilde)] :
    ∃ Cgen : ℝ,
      0 ≤ Cgen ∧
      1 ≤ Cgen ∧
      ∀ (Λ : Set ℕ) (t : ℕ → ℕ) (g : ℕ → α → ℂ) (N : ℝ)
        (f : α → ℂ)
        (x : WeakGridSpace.BesovishSpace
          (souzaAtomFamily G s p hs hp hp_top) q)
        (R : WeakGridSpace.LpGridRepresentation
          (souzaAtomFamily G s p hs hp hp_top)
          (x : Lp ℂ p G.toWeakGridSpace.measure)),
          0 ≤ N →
          WeakGridSpace.RepresentsFunction
            (G := G.toWeakGridSpace) (p := p) f
            (x : Lp ℂ p G.toWeakGridSpace.measure) →
          WeakGridSpace.LpGridRepresentation.FinitePQCost (q := q) R →
          (∀ i ∈ Λ,
            ∃ C : ℝ,
              SouzaPointwiseSelfsTailBound
                G β p qtilde hβ hp hp_top (t i) (g i) C) →
          (∀ k (Q : WeakGridSpace.LevelCell G.toWeakGridSpace k),
            (R.block k).coeff Q ≠ 0 →
              ∃ T : ℝ,
                HasSum
                  (fun i : {i // i ∈ Λ} =>
                    nonArchimedeanRelevantTailSelfsInfiniteTerm
                      G β p qtilde hβ hp hp_top Λ t g Q i)
                  T ∧
                T ≤ N) →
          (∀ k (Q : WeakGridSpace.LevelCell G.toWeakGridSpace k) i,
            i ∈ Λ →
              (R.block k).coeff Q ≠ 0 →
                goodGridLevelCellMeetsSupport G Q (g i) →
                  t i ≤ k) →
          ∃ h : α → ℂ,
            ∃ absSum : α → ℝ,
              (∀ᵐ z ∂G.toWeakGridSpace.measure,
                f z ≠ 0 →
                  HasSum
                    (fun i : {i // i ∈ Λ} => ‖g i.1 z‖)
                    (absSum z) ∧
                  absSum z ≤ Cgen * N) ∧
              (∀ᵐ z ∂G.toWeakGridSpace.measure,
                HasSum
                  (fun i : {i // i ∈ Λ} => g i.1 z * f z)
                  (h z)) ∧
              (∀ᵐ z ∂G.toWeakGridSpace.measure,
                ‖h z‖ ≤ Cgen * N * ‖f z‖) ∧
              (∃ hmem : MemLp h p G.toWeakGridSpace.measure,
                ‖MemLp.toLp h hmem‖ ≤
                  Cgen * N * ‖(x : Lp ℂ p G.toWeakGridSpace.measure)‖) ∧
              ∃ y : WeakGridSpace.BesovishSpace
                  (souzaAtomFamily G s p hs hp hp_top) q,
                ∃ S : WeakGridSpace.LpGridRepresentation
                    (souzaAtomFamily G s p hs hp hp_top)
                    (y : Lp ℂ p G.toWeakGridSpace.measure),
                  WeakGridSpace.RepresentsFunction
                    (G := G.toWeakGridSpace) (p := p) h
                    (y : Lp ℂ p G.toWeakGridSpace.measure) ∧
                  WeakGridSpace.LpGridRepresentation.FinitePQCost (q := q) S ∧
                  WeakGridSpace.LpGridRepresentation.pqCost (q := q) S ≤
                    Cgen * N *
                      WeakGridSpace.LpGridRepresentation.pqCost (q := q) R := ...
\end{lstlisting}

\begin{theorem}[Non-Archimedean property, positive version]\label{posrem}
In the setting of Theorems~\ref{sepa} and~\ref{sepainf} (finite or infinite
$\Lambda$, with the respective hypotheses, the tail seminorms being replaced
by their positive-cone counterparts
$$|g|_{\mathcal{B}^{\beta+,t}_{p,\tilde q,selfs}}
=\sup_{\substack{ a_Q \in \mathcal{A}^{sz}_{\beta,p}\\ Q\in \mathcal{P}^j,\ j\geq t}} |ga_Q|_{\mathcal{B}^{\beta+}_{p,\tilde q}}$$
for the functions $g_i$), assume that each $g_i$ is
$\mathcal{B}^{\beta,t_i}_{p,\tilde q}$-positive. Let $R$ be the given
canonical representation $(c_Q)_Q$ of $f$. Then the representation
$S=(d_Q)_Q$ of $(\sum_i g_i)f$ produced by
Theorems~\ref{sepa}/\ref{sepainf} can be chosen so that, in addition:
\begin{itemize}
\item[i.] ({\it positivity}) if $R$ is a positive representation
($c_Q \geq 0$ for all $Q$), then $S$ is a cone-positive representation: all
$d_Q$ are nonnegative reals and the corresponding Souza's atoms are a.e.\
nonnegative;
\item[ii.] ({\it support}) unconditionally --- without assuming $c_Q\geq 0$
--- every cell $Q$ active in $S$ ($d_Q \neq 0$) satisfies: there exists
$i \in \Lambda$ with $g_i \neq 0$ a.e.\ on $Q$.\footnote{Lean: \lean{souzaNonArchimedeanPropertyPositiveCone} (finite
$\Lambda$) and \lean{souzaNonArchimedeanPropertyPositiveConeInfinite}
($\Lambda \subset \mathbb{N}$ infinite, with condition A as an
$[0,\infty]$-valued series bound) in
\lean{GoodGrid/Multipliers/NonArchimedeanPropertyPositiveStandalone.lean},
with cores in \lean{NonArchimedeanProperty.lean} and
\lean{GoodGrid/PositiveCone.lean}; the auxiliary product representation is
\lean{exists\_nonArchimedeanProductRepresentation\_positive}. Compared with
paper Remark~18.5 the statement was made precise as follows: (a) the
hypothesis on the representation $R$ of $f$ is only that it is canonical
(arbitrary complex coefficients); (b) the two conclusions are separated
according to their true strengths --- the support conclusion ii.\ holds
unconditionally, and is weakened to its a.e.\ form ($g_i \neq 0$ a.e.\ on
$Q$ rather than $Q \subset \operatorname{supp} g_i$); (c) the positivity
conclusion i.\ produces a \emph{cone-positive} representation (nonnegative
real coefficients, a.e.\ nonnegative Souza's atoms), not necessarily one by
canonical atoms. Axioms of all four non-Archimedean theorems checked:
\lean{propext}, \lean{Classical.choice}, \lean{Quot.sound}.}
\end{itemize}
\end{theorem}
\begin{lstlisting}[style=lean,basicstyle=\ttfamily\scriptsize]
theorem souzaNonArchimedeanPropertyPositiveCone
    (G : GoodGridSpace (α := α))
    (s β : ℝ) (p q qtilde : ℝ≥0∞)
    (hs : 0 < s) (hβ : 0 < β) (hβs : s < β)
    (hβ_lt_inv : β < (p.toReal)⁻¹)
    (hp : 1 ≤ p) (hp_top : p ≠ ∞)
    [Fact (1 ≤ p)] [Fact (1 ≤ q)] [Fact (1 ≤ qtilde)] :
    ∃ Cgen : ℝ,
      0 ≤ Cgen ∧
      ∀ (Λ : Finset ℕ) (t : ℕ → ℕ) (g : ℕ → α → ℂ) (N : ℝ)
        (f : α → ℂ)
        (x : WeakGridSpace.BesovishSpace
          (souzaAtomFamily G s p hs hp hp_top) q)
        (R : WeakGridSpace.LpGridRepresentation
          (souzaAtomFamily G s p hs hp hp_top)
          (x : Lp ℂ p G.toWeakGridSpace.measure)),
          0 ≤ N →
          WeakGridSpace.RepresentsFunction
            (G := G.toWeakGridSpace) (p := p) f
            (x : Lp ℂ p G.toWeakGridSpace.measure) →
          WeakGridSpace.LpGridRepresentation.FinitePQCost (q := q) R →
          SouzaCanonicalRepresentation G s p hs hp hp_top R →
          (∀ i ∈ Λ,
            SouzaPositiveFunction G β p qtilde hβ hp hp_top (g i)) →
          (∀ i ∈ Λ,
            ∃ C : ℝ≥0∞,
              SouzaPositivePointwiseSelfsTailBound
                G β p qtilde hβ hp hp_top (t i) (g i) C) →
          (∀ k (Q : WeakGridSpace.LevelCell G.toWeakGridSpace k),
            (R.block k).coeff Q ≠ 0 →
              nonArchimedeanRelevantPositiveTailSelfsSum
                G β p qtilde hβ hp hp_top Λ t g Q ≤ ENNReal.ofReal N) →
          (∀ k (Q : WeakGridSpace.LevelCell G.toWeakGridSpace k) i,
            i ∈ Λ →
              (R.block k).coeff Q ≠ 0 →
                goodGridLevelCellMeetsSupport G Q (g i) →
                  t i ≤ k) →
          ∃ y : WeakGridSpace.BesovishSpace
              (souzaAtomFamily G s p hs hp hp_top) q,
            ∃ S : WeakGridSpace.LpGridRepresentation
                (souzaAtomFamily G s p hs hp hp_top)
                (y : Lp ℂ p G.toWeakGridSpace.measure),
              WeakGridSpace.RepresentsFunction
                (G := G.toWeakGridSpace) (p := p)
                (fun z => (∑ i ∈ Λ, g i z) * f z)
                (y : Lp ℂ p G.toWeakGridSpace.measure) ∧
              WeakGridSpace.LpGridRepresentation.FinitePQCost (q := q) S ∧
              WeakGridSpace.LpGridRepresentation.pqCost (q := q) S ≤
                Cgen * N *
                  WeakGridSpace.LpGridRepresentation.pqCost (q := q) R ∧
              -- [ii] support: holds with only canonical atoms on `R` (no sign on `c_Q`).
              (∀ k (Q : WeakGridSpace.LevelCell G.toWeakGridSpace k),
                (S.block k).coeff Q ≠ 0 →
                  ∃ i ∈ Λ, ∀ᵐ z ∂(G.toWeakGridSpace.measure.restrict Q.1), g i z ≠ 0) ∧
              -- [i] positivity: only when `R` is fully positive (canonical + `c_Q ≥ 0`).
              (SouzaPositiveRepresentation G s p hs hp hp_top R →
                SouzaConePositiveRepresentation G s p hs hp hp_top S) := ...

theorem souzaNonArchimedeanPropertyPositiveConeInfinite
    (G : GoodGridSpace (α := α))
    (s β : ℝ) (p q qtilde : ℝ≥0∞)
    (hs : 0 < s) (hβ : 0 < β) (hβs : s < β)
    (hβ_lt_inv : β < (p.toReal)⁻¹)
    (hp : 1 ≤ p) (hp_top : p ≠ ∞)
    [Fact (1 ≤ p)] [Fact (1 ≤ q)] [Fact (1 ≤ qtilde)] :
    ∃ Cgen : ℝ,
      0 ≤ Cgen ∧
      1 ≤ Cgen ∧
      ∀ (Λ : Set ℕ) (t : ℕ → ℕ) (g : ℕ → α → ℂ) (N : ℝ)
        (f : α → ℂ)
        (x : WeakGridSpace.BesovishSpace
          (souzaAtomFamily G s p hs hp hp_top) q)
        (R : WeakGridSpace.LpGridRepresentation
          (souzaAtomFamily G s p hs hp hp_top)
          (x : Lp ℂ p G.toWeakGridSpace.measure)),
          0 ≤ N →
          WeakGridSpace.RepresentsFunction
            (G := G.toWeakGridSpace) (p := p) f
            (x : Lp ℂ p G.toWeakGridSpace.measure) →
          WeakGridSpace.LpGridRepresentation.FinitePQCost (q := q) R →
          SouzaCanonicalRepresentation G s p hs hp hp_top R →
          (∀ i ∈ Λ,
            SouzaPositiveFunction G β p qtilde hβ hp hp_top (g i)) →
          (∀ i ∈ Λ,
            ∃ C : ℝ≥0∞,
              SouzaPositivePointwiseSelfsTailBound
                G β p qtilde hβ hp hp_top (t i) (g i) C) →
          (∀ k (Q : WeakGridSpace.LevelCell G.toWeakGridSpace k),
            (R.block k).coeff Q ≠ 0 →
              (∑' i : {i // i ∈ Λ},
                nonArchimedeanRelevantPositiveTailSelfsInfiniteTerm
                  G β p qtilde hβ hp hp_top Λ t g Q i) ≤ ENNReal.ofReal N) →
          (∀ k (Q : WeakGridSpace.LevelCell G.toWeakGridSpace k) i,
            i ∈ Λ →
              (R.block k).coeff Q ≠ 0 →
                goodGridLevelCellMeetsSupport G Q (g i) →
                  t i ≤ k) →
          ∃ h : α → ℂ,
            ∃ absSum : α → ℝ,
              (∀ᵐ z ∂G.toWeakGridSpace.measure,
                f z ≠ 0 →
                  HasSum
                    (fun i : {i // i ∈ Λ} => ‖g i.1 z‖)
                    (absSum z) ∧
                  absSum z ≤ Cgen * N) ∧
              (∀ᵐ z ∂G.toWeakGridSpace.measure,
                HasSum
                  (fun i : {i // i ∈ Λ} => g i.1 z * f z)
                  (h z)) ∧
              (∀ᵐ z ∂G.toWeakGridSpace.measure,
                ‖h z‖ ≤ Cgen * N * ‖f z‖) ∧
              (∃ hmem : MemLp h p G.toWeakGridSpace.measure,
                ‖MemLp.toLp h hmem‖ ≤
                  Cgen * N * ‖(x : Lp ℂ p G.toWeakGridSpace.measure)‖) ∧
              ∃ y : WeakGridSpace.BesovishSpace
                  (souzaAtomFamily G s p hs hp hp_top) q,
                ∃ S : WeakGridSpace.LpGridRepresentation
                    (souzaAtomFamily G s p hs hp hp_top)
                    (y : Lp ℂ p G.toWeakGridSpace.measure),
                  WeakGridSpace.RepresentsFunction
                    (G := G.toWeakGridSpace) (p := p) h
                    (y : Lp ℂ p G.toWeakGridSpace.measure) ∧
                  WeakGridSpace.LpGridRepresentation.FinitePQCost (q := q) S ∧
                  WeakGridSpace.LpGridRepresentation.pqCost (q := q) S ≤
                    Cgen * N *
                      WeakGridSpace.LpGridRepresentation.pqCost (q := q) R ∧
                  (∀ k (Q : WeakGridSpace.LevelCell G.toWeakGridSpace k),
                    (S.block k).coeff Q ≠ 0 →
                      ∃ i ∈ Λ,
                        ∀ᵐ z ∂(G.toWeakGridSpace.measure.restrict Q.1),
                          g i z ≠ 0) ∧
                  (SouzaPositiveRepresentation G s p hs hp hp_top R →
                    SouzaConePositiveRepresentation G s p hs hp hp_top S) := ...
\end{lstlisting}

\begin{remark}
Paper Corollary~18.6 ($\mathcal{B}^\beta_{p,\tilde{q},selfs} \subset
M(\mathcal{B}^s_{p,q})$ continuously, for $\beta>s$) follows from the
formalized Theorem~\ref{sepa} exactly as in the paper, but it is not itself
stated in the repository; it is therefore omitted from this document.
\end{remark}

\begin{thebibliography}{9}

\bibitem{paper} D. Smania,
\emph{Besov-ish spaces through atomic decomposition},
Analysis \& PDE \textbf{15} (2022), no.~1, 123--174.
LaTeX source: \texttt{latex/besovish-2021-07-19.tex} in this repository.

\bibitem{gw} M. Girardi and W. Sweldens,
\emph{A new class of unbalanced Haar wavelets that form an unconditional basis for $L_p$ on general measure spaces},
J. Fourier Anal. Appl. \textbf{3} (1997), 457--474.

\bibitem{souzao1} G. S. de Souza,
\emph{Spaces formed by special atoms},
PhD thesis, SUNY Albany, 1980.

\bibitem{repo} D. Smania,
\emph{BesovSpacesGoodGrid: a Lean 4 formalization},
git repository, 2026. Dependencies: \texttt{UnbalancedHaarWavelet},
\texttt{UnconditionalSchauderBasis}, \texttt{LaminarFamiliesMaximalBinaryTrees},
\texttt{Burkholder} (SmaniaD).

\end{thebibliography}

\end{document}
